\exercisegroup

\begin{enumerate}
\def\labelenumi{\arabic{enumi})}
\tightlist
\item
  Let \(S\) be an arbitrary infinite set.
\end{enumerate}

\begin{enumerate}
\def\labelenumi{\alph{enumi})}
\tightlist
\item
  Show that the smallest cardinal of any total set in the normed space \(\mathcal{B}(\mathrm{S})\) of bounded mappings of S in \(\mathbf{R}\) (I, p.~4, Example) is equal to \(2^{\mathrm{Card}(\mathrm{S})}\) (consider the set of characteristic functions of subsets of S and note that there exists an enumerable set everywhere dense in \(\mathbf{R}\) ). b) Show that the smallest cardinal of any total set in the normed space \(\ell^1 (\mathrm{S})\) (I, p.~4, Example) is equal to \(\mathrm{Card}(\mathrm{S})\) .
\end{enumerate}

\begin{enumerate}
\def\labelenumi{\arabic{enumi})}
\setcounter{enumi}{1}
\item
  In the product topological vector space \(E = R^{N}\) over the field R, denote by \(e_{n} (n \in \mathbf{N})\) the elements of the canonical basis of the direct sum \(\mathbf{R}^{(\mathbf{N})}\) . Write \(a_{0} = e_{0}\) , \(a_{n} = e_{0} + (1/n)e_{n}\) for \(n \geqslant 1\) . Show that, for every integer \(n \geqslant 0\) , the \(a_{i}\) such that \(0 \leqslant i \leqslant n\) form a topologically independent family in E, but that the infinite family \((a_{n})_{n \geqslant 0}\) is not topologically independent. If M is the closed vector subspace \(Ra_{0}\) , the classes \(\dot{a}_{n}\) of the \(a_{n}\) in E/M form a topologically independent family (for \(n \geqslant 1\) ), but the closed vector subspace N generated by the \(a_{n}\) , with index \(n \geqslant 1\) , in E contains M.
\item
  Let E be a topological vector space over R, and f a homomorphism of the additive group of E in R. Show that if there exists a neighbourhood of 0 in E in which f is bounded, then f is a continuous linear form in E. This is so in particular when f is semi-continuous (lower or upper).
\item
  Denote by K the field R with the absolute value \(p(\xi) = |\xi|^{1/2}\) . Let E be the vector space over K of the real valued regulated functions defined over I = \([0, 1]\) , continuous on the right everywhere and zero at the point 1; show that on E the mapping \(x \mapsto \|x\| = \int_{0}^{1} |x(t)|^{1/2} dt\) is a norm. Show that for every function \(x \geqslant 0\) in E, there exists in E two functions \(x_{1} \geqslant 0\) , \(x_{2} \geqslant 0\) such that \(x = \frac{1}{2}(x_{1} + x_{2})\) and \(\|x_{1}\| = \|x_{2}\| = \frac{1}{\sqrt{2}} \|x\|\) . Deduce that every continuous linear form on E is identically zero.
\item
  Let K be a Hausdorff topological division ring of which the topology is locally retrobounded (GT, III, § 6, exerc. 22). Extend prop. 2 of I, p.~12 and th. 1 of I, p.~13 to topological vector spaces over K; similarly extend th. 2 of I, p.~13 and prop. 3 of I, p.~14 when K is also complete.
\item
  Let K be the topological division ring obtained by transferring the usual topology of \(Q^{2}\) to the field \(\mathbf{Q}(\sqrt{2})\) by the mapping \((x, y) \mapsto x + y\sqrt{2}\) .
\end{enumerate}

\begin{enumerate}
\def\labelenumi{\alph{enumi})}
\item
  Let \(\mathbf{E}\) be the set \(\mathbf{Q}(\sqrt{2})\) with its vector space structure over \(\mathbf{K}\) and with the topology induced by that of \(\mathbf{R}\) . Show that \(\mathbf{E}\) is a Hausdorff topological vector space, of dimension 1 on \(\mathbf{K}\) , but that it is not isomorphic to \(\mathbf{K}_s\) .
\item
  Let \(\mathbf{F}\) be the topological vector space \(\mathbf{E} \times \mathbf{E}\) over \(\mathbf{K}\) ; in \(\mathbf{F}\) , the hyperplane \(\mathbf{E} \times \{0\}\) is closed but there is no continuous linear form \(f\) on \(\mathbf{E} \times \mathbf{E}\) such that this hyperplane is given by the equation \(f(x) = 0\) .
\end{enumerate}

\begin{enumerate}
\def\labelenumi{\arabic{enumi})}
\setcounter{enumi}{6}
\item
  Let K be a valued division ring which is non-discrete and non-complete, let E be the topological vector subspace \(K + Ka\) of \(\hat{K}\) where \(a \notin K\) , and let F be the product space \(K \times E\) . In F, the subspace \(M = K \times \{0\}\) is closed and of codimension 2. Let N be the complementary subspace to M in F generated by the vectors (0, 1) and (1, a); show that F is not the direct topological sum of M and N.
\item
  Let \(p\) be a prime number, \(\mathbf{Q}_p\) the field of \(p\) -adic numbers (GT, III, § 6, exerc. 23). Let \(\mathrm{E}_0\) be the topological product space \(\mathbf{Q}_p \times \mathbf{R}\) ; if \(\mathbf{K}\) denotes the field \(\mathbf{Q}\) with the discrete topology, then \(\mathrm{E}_0\) is a topological vector space over \(\mathbf{K}\) . Let \(\mathbf{M}\) be the vector subspace formed by the elements \((r, r)\) where \(r\) varies in \(\mathbf{Q}\) ; further let \(\theta\) be an irrational number and \(\mathbf{N}\) the vector subspace formed by the elements \((0, r\theta)\) where \(r\) varies in \(\mathbf{Q}\) . Let \(\mathbf{E}\) be the subspace \(\mathbf{M} + \mathbf{N}\) of \(\mathrm{E}_0\) ; show that \(\mathbf{N}\) is a closed hyperplane in \(\mathbf{E}\) , but that there does not exist a complementary topological subspace to \(\mathbf{N}\) (note that \(\mathbf{M}\) is everywhere dense in \(\mathrm{E}_0\) ).
\item
  Let \(\mathbf{X}\) be a Hausdorff topological space, and let \(\mathbf{V}\) be a vector subspace of finite dimension \(n\) of the space \(\mathcal{C}(\mathbf{X};\mathbf{R})\) .
\end{enumerate}

\begin{enumerate}
\def\labelenumi{\alph{enumi})}
\item
  Show that there exist \(n\) pair-wise disjoint, open sets \(\mathrm{U}_i\) ( \(1 \leqslant i \leqslant n\) ) in \(\mathbf{X}\) , such that any function \(f \in \mathbf{V}\) which is identically zero in each of the \(\mathrm{U}_i\) , is identically zero in \(\mathbf{X}\) (use A, II, § 7.5, cor. 3).
\item
  Let \(x_{i} \in \mathbf{U}_{i}\) for \(1 \leqslant i \leqslant n\) . Deduce from a) that there exists a constant \(c > 0\) such that, for every function \(f \in \mathbf{V}\) , we have
\end{enumerate}

\[
\sup _ {x \in \mathrm{X}} | f (x) | \leqslant c \sum_ {i = 1} ^ {n} | f (x _ {i}) |.
\]

\begin{enumerate}
\def\labelenumi{\arabic{enumi})}
\setcounter{enumi}{9}
\tightlist
\item
  Let \(\mathbf{K}\) be a locally compact non-discrete valued division ring, and \(\mathbf{E}\) a left vector space of finite dimension over \(\mathbf{K}\) . Denote by \(\mathfrak{N}(\mathbf{E})\) the set of norms on \(\mathbf{E}\) , which is a subspace of the space \(\mathcal{C}(\mathbf{E};\mathbf{R})\) of mappings of \(\mathbf{E}\) , continuous (in the canonical topology), in \(\mathbf{R}\) .
\end{enumerate}

\begin{enumerate}
\def\labelenumi{\alph{enumi})}
\item
  When we give to \(\mathcal{C}(\mathrm{E};\mathbf{R})\) the compact convergence topology \(*\) (for which it is a Fréchet space), the set \(\mathfrak{N}(\mathrm{E})\) is closed in \(\mathcal{C}(\mathrm{E};\mathbf{R})\) , and locally compact.
\item
  Let \(p_0\) be an element of \(\mathfrak{N}(\mathrm{E})\) ; show that there exists a continuous mapping \((\lambda, p) \mapsto \pi_{\lambda}(p)\) of \([0, 1] \times \mathfrak{N}(\mathrm{E})\) in \(\mathfrak{N}(\mathrm{E})\) such that \(\pi_0(p) = p\) and \(\pi_1(p) = p_0\) for every \(p \in \mathfrak{N}(\mathrm{E})\) .
\end{enumerate}

\begin{enumerate}
\def\labelenumi{\arabic{enumi})}
\setcounter{enumi}{10}
\tightlist
\item
  With the hypotheses of I, p.~23, exerc. 7 show that if K and E are complete then every closed subspace of E has a topological complement (proceed as in a), loc. cit.).
\end{enumerate}

¶ 12) Let K be a locally compact valued division ring whose absolute value is non-discrete and ultrametric. We call a norm on the left vector space E over K an ultranorm if it satisfies the ultrametric inequality (II, p.~2).

\begin{enumerate}
\def\labelenumi{\alph{enumi})}
\tightlist
\item
  Let E be a finite dimensional left vector space over K, let \(\alpha\) be an ultranorm on E and H a hyperplane in E given by the equation \(\langle x, a^{*} \rangle = 0\) . Show that there exists a point \(x_{0} \in E\) at which the function \(x \mapsto |\langle x, a^{*} \rangle| / \alpha(x)\) attains its upper bound in \(E \setminus \{0\}\) ; show that then
\end{enumerate}

\[
\alpha (x) = \sup \left(\alpha \left(x - \frac {\langle x , a ^ {*} \rangle}{\langle x _ {0} , a ^ {*} \rangle} x _ {0}\right), \frac {| \langle x , a ^ {*} \rangle |}{| \langle x _ {0} , a ^ {*} \rangle |} \alpha (x _ {0})\right).
\]

Deduce that there exists a basis \((a_{i})\) of E and a family \((r_i)\) of real numbers \(>0\) such that, for all \(x = \sum_{i} \xi_{i} a_{i}\) we have \(\alpha(x) = \sup_{i} (r_{i} |\xi_{i}|)\) . We say that \(\alpha\) is in the standard form relative to the basis \((a_{i})\) .

\begin{enumerate}
\def\labelenumi{\alph{enumi})}
\setcounter{enumi}{1}
\tightlist
\item
  Let \(\alpha^{*}\) be the norm on E* the dual of E canonically associated with \(\alpha\) by
\end{enumerate}

\[
\alpha^ {*} (x ^ {*}) = \sup _ {x \neq 0} | \langle x, x ^ {*} \rangle | / \alpha (x);
\]

it is an ultranorm. Show that for all \(x_0 \neq 0\) in \(\mathbf{E}\) , there exists \(x_0^* \in \mathbf{E}^*\) such that \(\alpha(x_0) = |\langle x_0, x_0^* \rangle| / \alpha^*(x_0^*)\) .

\begin{enumerate}
\def\labelenumi{\alph{enumi})}
\setcounter{enumi}{2}
\tightlist
\item
  Let \(\alpha, \beta\) be any two ultranorms on E. Show that there exists a basis of E such that relative to this basis \(\alpha\) and \(\beta\) are both of the standard form (consider a point \(x_0 \in \mathrm{E} \setminus \{0\}\) at which \(\alpha / \beta\) attains its maximum; then use \(b\) ) and proceed by induction on \(\dim \mathrm{E}\) ).
\end{enumerate}

\(d)\) Let \(\mathfrak{N}_0(\mathrm{E})\) , the set of ultranorms on \(\mathrm{E}\) , be considered as a subspace of \(\mathfrak{N}(\mathrm{E})\) (exerc. 10). Show that \(\mathfrak{N}_0(\mathrm{E})\) is closed in \(\mathfrak{N}(\mathrm{E})\) . Let \(\alpha_0\) be an element of \(\mathfrak{N}_0(\mathrm{E})\) ; for each \(\alpha \in \mathfrak{N}_0(\mathrm{E})\) and for \(0 \leqslant t \leqslant 1\) , let \(\mathrm{P}_{\alpha}(t)\) be the set of \(\beta \in \mathfrak{N}_0(\mathrm{E})\) such that \(\beta(x) \leqslant \alpha_0(x)^{1 - t}\alpha(x)^t\) for all \(x \in \mathrm{E}\) . Show that \(\mathrm{P}_{\alpha}(t)\) is not empty and that \(\pi_t^\alpha = \sup \mathrm{P}_{\alpha}(t)\) is an ultranorm. Further, the mapping \((t, \alpha) \mapsto \pi_t^\alpha\) of \([0, 1] \times \mathfrak{N}_0(E)\) in \(\mathfrak{N}_0(E)\) is continuous and such that \(\pi_0^\alpha = \alpha_0\) and \(\pi_1^\alpha = \alpha\) (use \(c\) ).

* e) Let A be the ring of the absolute value of K, m its maximal ideal such that \(k = A/m\) is a finite field with q elements (CA, VI, § 5, No.~1, prop. 2). For every ultranorm \(\alpha\) on E, the image \(X_{\alpha}\) of the set of values of \(\log \alpha(x)\) for \(x \in E \setminus \{0\}\) under the canonical mapping in the quotient group \(\mathbf{R}/(\mathbf{Z}, \log q)\) is a finite set having at most \(n = \dim E\) elements (use a); the number of these elements is denoted by \(r(\alpha)\) and called the rank of \(\alpha\) . Show that r is a lower semi-continuous mapping of \(\mathfrak{N}_{0}(E)\) in N and that the set \(\mathfrak{N}_{0}'(E)\) of the \(\alpha\) for which \(r(\alpha) = n\) is open and everywhere dense in \(\mathfrak{N}_{0}(E)\) (use a) and c)).

\begin{enumerate}
\def\labelenumi{\alph{enumi})}
\setcounter{enumi}{5}
\tightlist
\item
  Suppose that \(r(\alpha) = n\) ; let \((a_i)\) be a basis of E relative to which \(\alpha\) is of the standard form; show that there exists a neighbourhood V of \(\alpha\) in \(\mathfrak{N}_0'(\mathrm{E})\) such that every \(\beta \in V\) has the standard form relative to \((a_i)(\text{use } b)\) ; deduce that there is a neighbourhood \(W \subset V\) of \(\alpha\) homeomorphic to an open set in \(R^n\) .
\end{enumerate}

\(g)\) For every basis \((a_i)\) of E, show that the set of ultranorms \(\alpha\) that have a standard form relative to \((a_i)\) is closed in \(\mathfrak{N}_0^{\prime}(\mathrm{E})\) . Deduce that if \(\alpha \in \mathfrak{N}_0^{\prime}(\mathrm{E})\) has a standard form relative to \((a_i)\) the same is true of all elements of the connected component containing \(\alpha\) in \(\mathfrak{N}_0^{\prime}(\mathrm{E})\) .

¶ 13) * We keep the general hypotheses and the notations of exerc. 12.

\begin{enumerate}
\def\labelenumi{\alph{enumi})}
\item
  Let \(\mathbf{L}\) be a free sub-A-module of \(\mathbf{E}\) of dimension \(n = \dim \mathbf{E}\) . For all \(x \in \mathbf{E} \setminus \{0\}\) , the set of the \(a \in \mathbf{A}\) such that \(ax \in \mathbf{L}\) is a fractional ideal of \(\mathbf{K}\) of the form \(\mathfrak{m}^h\) ( \(h\) a positive or negative integer); putting \(\alpha(x) = q^h\) and \(\alpha(0) = 0\) , show that \(\alpha\) is an ultranorm on \(\mathbf{E}\) . It is said to be associated with the free A-module \(\mathbf{L}\) .
\item
  Conversely, if \(\alpha\) is an ultranorm on E, the set \(L_{\alpha}\) of the \(x \in E\) such that \(\alpha(x) \leqslant 1\) is a free A-module of dimension \(n\) . If \([\alpha]\) is the norm associated with \(L_{\alpha}\) , we have \(\alpha \leqslant [\alpha] \leqslant q\alpha\) , and \([\alpha]\) is the lower bound of the norms associated with free A-modules and which are \(\geqslant \alpha\) . We have \([q\alpha] = q \cdot [\alpha]\) , and \(\alpha(x) = \inf(q^{-t}[q^t\alpha](x))\) for all \(x \in E\) , where \(t\) varies in the interval {[}0, 1{]}. Further, the function \(t \mapsto [q^t\alpha](x)\) is left continuous in this interval.
\item
  With the same notations, show that for \(0 \leqslant t \leqslant 1\) , there are at most n distinct ultranorms among the \([q^{t}\alpha]\) . Conversely, let L be the set of ultranorms associated with the free A-modules of dimension n, and let \((\alpha_{t})_{0 \leqslant t \leqslant 1}\) be an increasing family of ultranorms of L such that \(\alpha_{1} = q\alpha_{0}\) . Show that there exists a basis of E relative to which all the \(\alpha_{t}\) have the standard form (if \(u \in A\) is an element of valuation 1, and \(L_{t}\) the free A-module of the \(x \in E\) such that \(\alpha_{t}(x) \leqslant 1\) , consider the vector spaces \(L_{t}/uL_{0}\) on k). Deduce further, that if, for all \(x \in E\) , \(t \mapsto \alpha_{t}(x)\) is left-continuous in \([0, 1]\) then there exists a unique ultranorm \(\alpha\) such that \(\alpha_{t} = [q^{t}\alpha]\) for all \(t \in [0, 1]\) .
\end{enumerate}

\(d\) ) The linear group \(\mathbf{GL}(\mathrm{E})\) operates continuously in \(\mathfrak{N}_0(\mathrm{E})\) ; show that it operates properly. For all \(\alpha \in \mathfrak{N}_0(\mathrm{E})\) , the stabiliser \(S_{\alpha}\) of \(\alpha\) in \(\mathbf{GL}(\mathrm{E})\) is the intersection of the stabilisers of the \([q^t\alpha]\) for \(0\leqslant t\leqslant 1\) ; deduce that \(S_{\alpha}\) is an open compact subgroup of \(\mathbf{GL}(\mathrm{E})\) , and hence that the orbit of each \(\alpha \in \mathfrak{N}_0(\mathrm{E})\) is a closed, discrete subspace of \(\mathfrak{N}_0(\mathrm{E})\) .

\begin{enumerate}
\def\labelenumi{\alph{enumi})}
\setcounter{enumi}{4}
\tightlist
\item
  For every ultranorm \(\alpha \in \mathfrak{N}_0(\mathrm{E})\) consider the decreasing sequence of the dimensions of the vector \(k\) -spaces \(\mathrm{L}_t / u\mathrm{L}_0\) , where \(\mathrm{L}_t\) is the A-module of the \(x \in \mathrm{E}\) such that \([q^t\alpha](x) \leqslant 1\) , and \(t\) varies from 0 to 1; we call this sequence, the sequence of invariants of \(\alpha\) . In order that \(\alpha\) and \(\beta\) belong to the same orbit in \(\mathfrak{N}_0(\mathrm{E})\) , it is necessary and sufficient that \(\mathrm{X}_{\alpha} = \mathrm{X}_{\beta}\) (exerc. 12, e)) and that the sequence of the invariants of \(\alpha\) and of \(\beta\) should be the same (use exerc. 12, b)). f) Deduce from e) that the space of the orbits \(\mathfrak{N}_0(\mathrm{E}) / \mathbf{GL}(\mathrm{E})\) is isomorphic with the space of the orbits \(\mathbf{T}^n / \mathfrak{S}_n\) , where the symmetric group operates on the right on \(\mathbf{T}^n\) by \((z_1, ..., z_n) \mapsto (z_{\sigma(1)}, ..., z_{\sigma(n)})\) . *
\end{enumerate}

\begin{enumerate}
\def\labelenumi{\arabic{enumi})}
\setcounter{enumi}{13}
\item
  Generalize the results of No.~2 and No.~3 to topological vector spaces E over a discrete division ring K, such that there exists a fundamental system of balanced neighbourhoods of 0 in E (i.e.~of neighbourhoods V such that K.V = V).
\item
  Let \(\mathbf{E}\) be a normed space of finite dimension \(n\) over \(\mathbf{R}\) or \(\mathbf{C}\) . Ascribe to the dual \(\mathbf{E}^*\) , the norm defined by \(\| x^* \| = \sup_{\| x \| \leqslant 1} |\langle x, x^* \rangle|\) (GT, X, § 3.2). Show that there exists a basis \((e_i)\)
\end{enumerate}

of E such that, if \((e_{i}^{*})\) is the dual basis, we have \(\|e_{i}\| = \|e_{i}^{*}\| = 1\) for all i. (Let \((a_{i})\) be a basis of E formed of vectors of norm 1; consider, the determinant \(\det(\xi_{ij})\) for each system of n vectors \(x_{i} = \sum_{j} \xi_{ij} a_{j}\) of norm 1, and consider such a system for which the absolute value of this determinant is maximal.)

\exercisegroup

\begin{enumerate}
\def\labelenumi{\arabic{enumi})}
\tightlist
\item
  \begin{enumerate}
  \def\labelenumii{\alph{enumii})}
  \tightlist
  \item
    Show that, if a Hausdorff topological vector space E over a non-discrete valued division ring K is such that every neighbourhood of 0 contains a vector subspace that is not the single point 0, then the topology of E cannot be defined by a norm. In particular, a product of an infinite sequence \((\mathrm{E}_{n})\) of Hausdorff topological vector spaces on K, none consisting of the single point 0, has a topology that cannot be defined by a norm.
  \end{enumerate}
\end{enumerate}

\begin{enumerate}
\def\labelenumi{\alph{enumi})}
\setcounter{enumi}{1}
\tightlist
\item
  Consider the product vector space \(E = K_{s}^{N}\) ; for all \(x = (\xi_{n}) \in E\) , put
\end{enumerate}

\[
| x | = \sum_ {n = 0} ^ {\infty} 2 ^ {- n} \left| \xi_ {n} \right| / (1 + | \xi_ {n} |).
\]

Show that the topology of E is defined by the distance \(d(x, y) = |x - y|\) , that \(|\lambda x| \leqslant |x|\) if \(|\lambda| \leqslant 1\) , \(|\lambda x| \leqslant |\lambda| \cdot |x|\) if \(|\lambda| \geqslant 1\) and that, for all \(x_0 \in \mathrm{E}\) , \(|\lambda x_0|\) tends to 0 with \(|\lambda|\) .

\begin{enumerate}
\def\labelenumi{\arabic{enumi})}
\setcounter{enumi}{1}
\tightlist
\item
  Let E and F be two complete, metrisable vector spaces over a non-discrete valued division ring, and let \(T_{0}\) be the topology of F. Let T be a Hausdorff topology on F, coarser than \(T_{0}\) . Show that if the linear mapping u of E in F is continuous for the topology T on F, it is still continuous for the topology \(T_{0}\) on F (use the cor. 5 of I, p.~19).
\end{enumerate}

Deduce that if \(T_{1}\) and \(T_{2}\) are two distinct topologies on a vector space E over a non-discrete valued division ring, compatible with the vector space structure of E, and for each of them E is metrisable and complete, then there does not exist a Hausdorff topology on E coarser than \(T_{1}\) and \(T_{2}\) . Give an example of two such topologies on an infinite dimensional vector space E (note that there exist bijections of E on itself such that both the bijection and its inverse are not continuous for a normed space topology on E).

\begin{enumerate}
\def\labelenumi{\arabic{enumi})}
\setcounter{enumi}{2}
\item
  Let E and F be two Hausdorff topological vector spaces over a non-discrete valued division ring; and suppose that E is metrisable and complete. Let u be a continuous linear injection of E in F, and let G be a vector subspace of \(u(\mathrm{E})\) ; suppose that there exists on G a topology T which is finer than the topology induced by that of F, is compatible with the vector space structure of G and for which G is metrisable and complete. Show that the mapping inverse to u, restricted to G, is continuous for T (use I, p.~19, cor. 5).
\item
  Let E, F be two complete metrisable vector spaces over a non-discrete valued division ring and let u be a continuous linear mapping of E in F. Show that if there exists in F a closed complementary subspace to \(u(\mathrm{E})\) , then \(u(\mathrm{E})\) is closed in F (use I, p.~19, cor. 5).
\item
  Let E and F be two complete metrisable vector spaces over a non-discrete valued division ring and let u be a linear mapping of E in F. Let N be the set of cluster points of u in F with respect to the filter of neighbourhoods of 0 in E; show that N is a closed vector subspace of F, and that, in order that u be continuous, it is necessary and sufficient that N be the single point 0 (use I, p.~19, cor. 5). Show that N is the smallest of the closed vector subspaces M of F such that, if \(\phi\) denotes the canonical homomorphism of F on F/M, then \(\phi \circ u\) is a continuous mapping of E in F/M.
\item
  Let E be a complete metrisable vector space over a non-discrete valued division ring K.
\end{enumerate}

\begin{enumerate}
\def\labelenumi{\alph{enumi})}
\tightlist
\item
  Let p be a lower semi-continuous mapping of E in the interval \([0, +\infty)\) of \(\overline{R}\) such that \(p(\lambda x) = |\lambda| \cdot p(x)\) for \(\lambda \neq 0\) in K and \(x \in E\) , such that \(p(0) = 0\) and \(p(x + y) \leqslant p(x) + p(y)\)
\end{enumerate}

for any \(x, y\) in E. Show that if \(p\) is finite in E, then \(p\) is continuous (consider the closed set B of \(x \in \mathbf{E}\) such that \(p(x) \leqslant 1\) , and use Baire's theorem).

\begin{enumerate}
\def\labelenumi{\alph{enumi})}
\setcounter{enumi}{1}
\tightlist
\item
  Let \((p_n)\) be a sequence of mappings of \(E\) in \([0, +\infty]\) satisfying the conditions of \(a\) ). Show that if none of the \(p_n\) are finite in \(E\) then there exists a point \(x \in E\) such that \(p_n(x) = +\infty\) for all \(n\) (same method as above).
\end{enumerate}

\begin{enumerate}
\def\labelenumi{\arabic{enumi})}
\setcounter{enumi}{6}
\tightlist
\item
  Let E be a complete metrisable vector space over a non-discrete valued division ring K. We say that a vector subspace M of E is paracomplete if there exists on M a complete metrisable vector space structure for which the canonical injection of M in E is continuous. a) Let M, N be two paracomplete subspaces of E such that \(M + N\) and \(M \cap N\) are closed in E. Show that M and N are closed in E. (Taking quotients by \(M \cap N\) reduces the question to the case where \(M \cap N = \{0\}\) , and we can consider then the mapping \((x, y) \mapsto x + y\) of \(M \times N\) in E).
\end{enumerate}

\begin{enumerate}
\def\labelenumi{\alph{enumi})}
\setcounter{enumi}{1}
\tightlist
\item
  Show that if E is the union of an increasing sequence of paracomplete subspaces, \((\mathbf{M}_{j})_{j\geq0}\) , then there exists an index j such that \(M_{j}=E\) . (Use Baire's theorem (GT, IX, § 5.3, th. 1) and I, p.~17, th. 1).
\end{enumerate}

\begin{enumerate}
\def\labelenumi{\arabic{enumi})}
\setcounter{enumi}{7}
\tightlist
\item
  Let E be a Banach space over a non-discrete valued division ring K. We say that a vector subspace M of E is strongly paracomplete if there exists a norm \(\|x\|_{M}\) on M for which M is a Banach space and the canonical injection of M in E is continuous.
\end{enumerate}

\begin{enumerate}
\def\labelenumi{\alph{enumi})}
\item
  Show that if \(\mathbf{M}\) and \(\mathbf{N}\) are two strongly paracomplete subspaces of \(\mathbf{E}\) , then \(\mathbf{M} + \mathbf{N}\) and \(\mathbf{M} \cap \mathbf{N}\) are also strongly paracomplete subspaces. (On \(\mathbf{M} + \mathbf{N}\) , consider the norm \(\|x\|_{\mathbf{M} + \mathbf{N}} = \inf (\|u\|_{\mathbf{M}} + \|v\|_{\mathbf{N}})\) , where the lower bound is taken over all pairs \((u, v)\) such that \(x = u + v\) , \(u \in \mathbf{M}\) and \(v \in \mathbf{N}\) .)
\item
  Let M, N be two strongly paracomplete subspaces of E such that N and \(M + N\) are closed. Show that \(\overline{M} = M + (\overline{M} \cap \overline{N})\) and \(\overline{M} \cap \overline{N} = \overline{M} \cap N\) (use exerc. 7, a)).
\end{enumerate}

\begin{enumerate}
\def\labelenumi{\arabic{enumi})}
\setcounter{enumi}{8}
\tightlist
\item
  \begin{enumerate}
  \def\labelenumii{\alph{enumii})}
  \tightlist
  \item
    Let \(a, b\) be two points of a normed space \(\mathbf{E}\) on the field \(\mathbf{R}\) . Denote by \(\delta(\mathbf{A})\) the diameter of a bounded set \(\mathbf{A}\) in \(\mathbf{E}\) (using the norm metric on \(\mathbf{E}\) ) and define inductively the sequence \((\mathbf{B}_n)_{n \geqslant 1}\) of bounded sets in \(\mathbf{E}\) satisfying the following conditions: \(\mathbf{B}_1\) is the set of those \(x \in \mathbf{E}\) such that \(\|x - a\| = \|x - b\| = \frac{1}{2}\|a - b\|\) ; for \(n > 1\) , \(\mathbf{B}_n\) is the set of those \(x \in \mathbf{B}_{n-1}\) such that \(\|x - y\| \leqslant \frac{1}{2}\delta(\mathbf{B}_{n-1})\) for all \(y \in \mathbf{B}_{n-1}\) . Show that the intersection of the \(\mathbf{B}_n\) is just the single point \(\frac{1}{2}(a + b)\) (note that \(\delta(\mathbf{B}_n) \leqslant \frac{1}{2}\delta(\mathbf{B}_{n-1})\) ).
  \end{enumerate}
\end{enumerate}

\begin{enumerate}
\def\labelenumi{\alph{enumi})}
\setcounter{enumi}{1}
\tightlist
\item
  Deduce from \(a)\) that if \(u\) is an isometry of the real Banach space \(E\) on the real Banach space \(F\) , then \(u\) is an affine linear mapping of \(E\) on \(F\) .
\end{enumerate}

\chapter{Convex sets and locally convex spaces}

In §§ 2 to 7 of this chapter, we shall be concerned only with vector spaces and affine spaces over the field of real numbers R, and when we speak of a vector space or an affine space without giving its division ring of scalars explicitly, then it is to be understood that this division ring is the field R. For vector spaces on C, see § 8.

\section{SEMI-NORMS}

Throughout this paragraph, K denotes a non-discrete valued division ring.

\subsection{Definition of semi-norms}

DEFINITION 1. --- Let E be a left vector space over K. A mapping p of E in \(R_{+} = \left[0, +\infty\right]\) , is called a semi-norm on E if it satisfies the following axioms:

\((\mathrm{SN}_{\mathrm{I}})\) If \(x \in \mathrm{E}\) and \(\lambda \in \mathrm{K}\) then \(p(\lambda x) = |\lambda| p(x)\) .

\((\mathrm{SN}_{\mathrm{II}})\) If \(x \in \mathrm{E}\) and \(y \in \mathrm{E}\) then \(p(x + y) \leqslant p(x) + p(y)\) .

Since \(p(x) \leqslant p(y) + p(x - y)\) and \(p(y) \leqslant p(x) + p(y - x)\) , from \(p(y - x) = p(x - y)\) , we deduce

\[
\left| p (x) - p (y) \right| \leqslant p (x - y).\tag{1}
\]

Examples. --- 1) A norm on E is a semi-norm p such that the relation \(p(x) = 0\) implies that x = 0 (I, p.~3).

\begin{enumerate}
\def\labelenumi{\arabic{enumi})}
\setcounter{enumi}{1}
\item
  For every linear form \(f\) on \(\mathbf{E}\) , the function \(x \mapsto |f(x)|\) is a semi-norm on \(\mathbf{E}\) .
\item
  If \(p_i (1 \leqslant i \leqslant n)\) is a finite set of semi-norms on \(E\) , then clearly \(p'(x) = \sup_{1 \leqslant i \leqslant n} p_i(x)\) .
\end{enumerate}

and \(p''(x) = \sum_{i=1}^{n} \alpha_i p_i(x)\) (where the \(\alpha_i\) are \(\geqslant 0\) ) are both semi-norms on E.

A mapping \(p\) of \(\mathbf{E}\) in \(\mathbf{R}_{+}\) is called an ultra-semi-norm if it satisfies \((\mathrm{SN}_{\mathrm{I}})\) and the following axiom:

\((\mathrm{SN}_{\mathrm{II}}^{\prime})\) If \(x \in \mathrm{E}\) and \(y \in \mathrm{E}\) , then \(p(x + y) \leqslant \sup(p(x), p(y))\) .

Clearly an ultra-semi-norm is a semi-norm.

To say that the absolute value on K is ultrametric (CA, VI, § 6.2) means that it is an ultra-semi-norm on the left vector space \(K_{s}\) , which is not identically zero.

PROPOSITION 1. --- Let E be a left topological vector space over K and let p be a semi-norm on E. The following conditions are equivalent :

\begin{enumerate}
\def\labelenumi{\alph{enumi})}
\item
  \(p\) is continuous in E.
\item
  \(p\) is continuous at the point 0.
\item
  \(p\) is uniformly continuous.
\item
  For each real number \(\alpha > 0\) , the set \(\mathbf{W}(p, \alpha)\) , of those \(x \in \mathbf{E}\) for which \(p(x) < \alpha\) , is open in \(\mathbf{E}\) .
\item
  There exists a real number \(\alpha > 0\) , such that \(\mathrm{W}(p, \alpha)\) is a neighbourhood of 0 in E.
\item
  For every real number \(\alpha > 0\) , the set \(\mathrm{V}(p, \alpha)\) , of those \(x \in \mathrm{E}\) for which \(p(x) \leqslant \alpha\) , is a neighbourhood of 0 in \(\mathrm{E}\) .
\end{enumerate}

In fact, the implications \(c) \Rightarrow a) \Rightarrow b) \Rightarrow d) \Rightarrow e) \Rightarrow f) \Rightarrow c)\) follow immediately from \((\mathrm{SN}_{\mathrm{I}})\) and inequality (1).

COROLLARY. --- If p is a continuous semi-norm on E and q is a semi-norm such that \(q \leqslant p\) , then q is continuous in E.

When \(p\) is an ultra-semi-norm on \(\mathbf{E}\) , then the sets \(\mathrm{W}(p, \alpha)\) and \(\mathrm{V}(p, \alpha)\) are both open and closed. For, we have seen that \(\mathrm{W}(p, \alpha)\) is open; on the other hand if \(z\) is a cluster point of \(\mathrm{W}(p, \alpha)\) , then there exists \(y \in \mathrm{W}(p, \alpha)\) such that \(p(y - z) < \alpha\) , and from \((\mathrm{SN}_{\mathrm{II}}')\) we have \(p(z) < \alpha\) , thus \(\mathrm{W}(p, \alpha)\) is closed. Also, \(\mathrm{V}(p, \alpha)\) is closed since \(p\) is continuous; further if \(p(x) \leqslant \alpha\) and \(p(y) \leqslant \alpha\) , then \(p(x + y) \leqslant \alpha\) by \((\mathrm{SN}_{\mathrm{II}}')\) , and this shows that \(\mathrm{V}(p, \alpha)\) is open.

\subsection{Topologies defined by semi-norms}

Let \(p\) be a semi-norm on the vector space \(E\) over \(K\) ; for every \(\alpha > 0\) let \(V(p, \alpha)\) be the subset of those \(x\) of \(E\) for which \(p(x) \leqslant \alpha\) . Clearly, if \(x \in V(p, \alpha)\) and \(\lambda \in K\) is such that \(|\lambda| \leqslant 1\) , then \(\lambda x \in V(p, \alpha)\) , in other words \(V(p, \alpha)\) is balanced. Further, for every \(x_0 \in E\) , there exists a non-zero scalar \(\mu \in K\) such that \(|\mu| \geqslant p(x_0) \alpha^{-1}\) , therefore \(\mu^{-1}x_0 \in V(p, \alpha)\) that is to say \(V(p, \alpha)\) is absorbent. Finally, from \((\mathrm{SN}_{\mathrm{II}})\) , we have \(V(p, \alpha/2) + V(p, \alpha/2) \subset V(p, \alpha)\) , and from \((\mathrm{SN}_{\mathrm{I}})\) that for every non-zero scalar \(\lambda\) in \(K\) we have \(\lambda V(p, \alpha) = V(p, |\lambda| \alpha)\) . We conclude from these remarks, by I, p.~7, prop. 4, that, when \(\alpha\) varies in the set of numbers \(> 0\) (or only in a sequence of strictly positive numbers tending to 0) then the sets \(V(p, \alpha)\) constitute a fundamental system of neighbourhoods of 0 for a topology compatible with the vector space structure of \(E\) ; we say that this topology is defined by the semi-norm \(p\) . A vector space \(E\) with such a topology is called a semi-normed space. Note that if \(W(p, \alpha)\) is the subset of \(x\) of \(E\) such that \(p(x) < \alpha\) , then the \(W(p, \alpha)\) constitute (where \(\alpha > 0\) , or \(\alpha\) )

varies in a strictly positive sequence of numbers tending to zero) a fundamental system of neighbourhoods of 0 for the topology defined by \(p\) .

If \(\Gamma\) is a set of semi-norms on \(\mathbf{E}\) , then the upper bound of the topologies defined by the semi-norms \(p \in \Gamma\) is compatible with the vector space structure (I, p.~10, cor. 4). A fundamental system of neighbourhoods of 0, for this topology, is given by the finite intersections \(\bigcap_{i} \mathrm{V}(p_i, \alpha_i)\) where \(p_i \in \Gamma\) and \(\alpha_i > 0\) . This topology is said to be defined

by the set of semi-norms \(\Gamma\) . It is the coarsest topology on \(\mathbf{E}\) amongst those that are invariant under all translations and for which the semi-norms \(p \in \Gamma\) are continuous.

Let E be a topological vector space over K : a system of semi-norms on E, say \(\Gamma\) , is called a fundamental system of semi-norms if the topology on E is the same as the topology defined by \(\Gamma\) .

Let E be a vector space over K, with the topology defined by a set of semi-norms \(\Gamma\) . For every semi-norm p, we have \(p(x - z) \leqslant p(x - y) + p(y - z)\) , which shows that the function \((x, y) \mapsto p(x - y)\) is a pseudometric on E (GT, IX, §1.1): it follows from the definitions that, when p varies in \(\Gamma\) , the set of these pseudometrics defines the uniform structure of the topological vector space E.

Remarks. --- 1) The topology defined by a finite set of semi-norms \(p_i\) ( \(1 \leqslant i \leqslant n\) ) on E, can be defined by the single semi-norm \(p = \sup_{1 \leqslant i \leqslant n} p_i\) . But a topology defined by

an infinite set of semi-norms cannot, in general, be defined by a single semi-norm (III, p.~37, exerc. 2).

\begin{enumerate}
\def\labelenumi{\arabic{enumi})}
\setcounter{enumi}{1}
\item
  Let \((\mathcal{T}_{\iota})_{\iota\in I}\) be a family of topologies on a vector space E over K, each of which is defined by a family of semi-norms \(\Gamma_{\iota}\) . Then the topology defined by the set of semi-norms \(\Gamma = \bigcup \Gamma_{\iota}\) is the upper bound of the topologies \(T_{\iota}\) .
\item
  If \(\Gamma_0\) is a set of semi-norms directed by the increasing order relation defined between two semi-norms \(p, q\) on E by « there exists \(\lambda > 0\) such that \(p \leqslant \lambda q\) », then a fundamental system of neighbourhoods of 0, for the topology defined by \(\Gamma_0\) , is obtained by taking the sets \(\mathbf{V}(p, \alpha)\) where \(p \in \Gamma_0\) and \(\alpha > 0\) . If \(\Gamma\) is any set of semi-norms on E, then a filtered set of semi-norms, defining the same topology as \(\Gamma\) , is the set \(\Gamma_0\) of upper envelopes of all finite families of semi-norms belonging to \(\Gamma\) .
\item
  Even if \(\mathbf{K} = \mathbf{R}\) , the topology of a topological vector space over \(\mathbf{K}\) cannot always be defined by a set of semi-norms (cf.~II, p.~24, corollary).
\end{enumerate}

Example. --- Let \(\mathcal{C}^{\infty}(\mathbf{R})\) be the vector space over R of real valued functions that are infinitely differentiable in R. For every function and every pair of integers \(n \geqslant 0\) , \(m \geqslant 1\) , put

\[
p _ {n, m} (f) = \sup _ {- m \leqslant t \leqslant m} | f ^ {(n)} (t) |\tag{2}
\]

with \(f^{(0)} = f\) . Obviously the \(p_{n,m}\) are semi-norms on \(\mathcal{C}^{\infty}(\mathbf{R})\) . In order that the functions \(f_{\alpha}\) tend to 0 (following a filter \(\mathfrak{F}\) on the set of indices) in \(\mathcal{C}^{\infty}(\mathbf{R})\) for the topology \(\mathcal{T}\) defined by the semi-norms \(p_{n,m}\) , it is necessary and sufficient that for all integers \(n \geqslant 0\) , the functions \(f_{\alpha}^{(n)}\) tend to 0 (following \(\mathfrak{F}\) ) uniformly on every compact subset of \(\mathbf{R}\) . We say that \(\mathcal{T}\) is the topology of compact convergence for the functions \(f \in \mathcal{C}^{\infty}(\mathbf{R})\) and all their derivatives (cf.~III, p.~9).

PROPOSITION 2. --- On a vector space E, let T be the topology defined by a set of semi-norms Γ.

\begin{enumerate}
\def\labelenumi{(\roman{enumi})}
\item
  The closure of \(\{0\}\) in \(\mathbf{E}\) , for \(\mathcal{T}\) , is the subset of \(x \in \mathbf{E}\) for which \(p(x) = 0\) for every semi-norm \(p \in \Gamma\) .
\item
  If \(\mathcal{T}\) is Hausdorff and \(\Gamma\) is enumerable, then \(\mathcal{T}\) is metrisable.
\end{enumerate}

The proposition follows immediately from the definitions and from GT, IX, § 2.4, cor. 1.

Note that if \(\mathcal{T}\) is metrisable, it may be that \(\mathcal{T}\) cannot be defined by a single norm; this is the case in the example given above (cf.~IV, p.~18, Example 4).

Let E be a vector space over K, with the topology defined by a set of semi-norms \(\Gamma\) . Let \(\hat{E}\) be the Hausdorff completion of E (I, p.~6), and \(\hat{\Gamma}\) be the set of mappings \(\hat{p}\) of \(\hat{E}\) in \(R_{+}\) where p varies in \(\Gamma\) (GT, II, § 3.7, prop. 15). By the principle of extending inequalities, the functions \(\hat{p} \in \hat{\Gamma}\) are semi-norms on \(\hat{E}\) , and the functions \(\hat{p}(x - y)\) form a set of pseudometrics defining the uniform structure of \(\hat{E}\) (GT, IX, § 1.3, prop. 1). We see, therefore, that \(\hat{\Gamma}\) is a fundamental set of semi-norms defining the topology of \(\hat{E}\) .

\subsection{Semi-norms in quotient spaces and in product spaces}

Let E be a topological vector space over K, whose topology is defined by \(\Gamma\) , a set of semi-norms. Clearly, the restrictions of the semi-norms of \(\Gamma\) to a vector sub-space M of E, define the topology induced on M by that of E.

Let \(\phi\) be the canonical mapping of E on the vector quotient space E/M. We show that, for every semi-norm \(p\) on E, the function

\[
\dot {p} (z) = \inf _ {\phi (x) = z} p (x)\tag{3}
\]

is a semi-norm on E/M. In fact, it is clear that \(\dot{p}\) satisfies the condition \((\mathrm{SN}_{\mathrm{I}})\) ; on the other hand, if \(z'\) , \(z''\) are two vectors of E/M, we have :

\[
\begin{array}{r l} \inf _ {\phi (x) = z ^ {\prime} + z ^ {\prime \prime}} p (x) & \leqslant \inf _ {\phi (x ^ {\prime}) = z ^ {\prime}, \phi (x ^ {\prime \prime}) = z ^ {\prime \prime}} p (x ^ {\prime} + x ^ {\prime \prime}) \\ & \leqslant \inf _ {\phi (x ^ {\prime}) = z ^ {\prime}, \phi (x ^ {\prime \prime}) = z ^ {\prime \prime}} \left(p (x ^ {\prime}) + p (x ^ {\prime \prime})\right) \\ & = \inf _ {\phi (x ^ {\prime}) = z ^ {\prime}} p (x ^ {\prime}) + \inf _ {\phi (x ^ {\prime \prime}) = z ^ {\prime \prime}} p (x ^ {\prime \prime}) \end{array}
\]

which shows that \(\dot{p}\) verifies \((\mathrm{SN}_{\mathrm{II}})\) . We say that \(\dot{p}\) is the quotient semi-norm of p by M.

The same reasoning proves that, if p is an ultra-semi-norm, then so also is \(\dot{p}\) .

This being so, we have (in the notation of No.~2)

\[
\phi (\mathrm{W} (p, \alpha)) = \mathrm{W} (\dot {p}, \alpha).\tag{4}
\]

for every \(\alpha > 0\) . In fact, to say that \(\dot{p}(z) < \alpha\) , means that there exists \(x \in \mathbf{E}\) such that \(\phi(x) = z\) and \(p(x) < \alpha\) , from which the relation (4) follows.

We deduce from this, that, if the set of semi-norms \(\Gamma\) is directed(II, p.~3, Remark 3), then the quotient topology on E/M is defined by the set of semi-norms \(\dot{p}\) , when p varies in \(\Gamma\) .

If N is the closure of \(\{0\}\) in E, the topology of E/N is defined by the quotient seminorms \(\dot{p}\) , where p varies in \(\Gamma\) (even if \(\Gamma\) is not filtered): here \(\dot{p}(\dot{x}) = p(x)\) for every x belonging to the class \(\dot{x} \mod N\) . Note that E/N is none other than the Hausdorff space associated with E (I, p.~4).

Let \(\mathbf{E}\) be a vector space over \(\mathbf{K}\) and \((\mathrm{E}_{\iota})_{\iota \in \mathbf{I}}\) be a family of vector spaces over \(\mathbf{K}\) , where \(\mathrm{E}_{\iota}\) has the topology \(\mathcal{T}_{\iota}\) defined by a set of semi-norms \(\Gamma_{\iota}\) . For each \(\iota \in \mathbf{I}\) , let \(f_{\iota}\) be a linear mapping of \(\mathbf{E}\) in \(\mathrm{E}_{\iota}\) ; clearly when \(p_{\iota}\) varies in the set \(\Gamma_{\iota}\) , then the \(p_{\iota} \circ f_{\iota}\) form a set \(\Gamma_{\iota}'\) of semi-norms on \(\mathbf{E}\) . The topology \(\mathcal{T}\) on \(\mathbf{E}\) , defined as being the coarsest of all those which make all the mappings \(f_{\iota}\) continuous (I, p.~9) is then defined by the set of semi-norms \(\Gamma' = \bigcup_{\iota \in \mathbf{I}} \Gamma_{\iota}'\) , this follows from the definition of neighbourhoods of 0 for \(\mathcal{T}\) (GT, I, § 2.3, prop. 4).

If the \(p_{\iota}\) are ultra-semi-norms, then so are the \(p_{\iota} \circ f_{\iota}\) .

Let E be a vector space over K, with the topology T defined by a family of semi-norms \((p_{\iota})_{\iota \in I}\) ; for every \(\iota \in I\) , let \(T_{\iota}\) be the topology defined by the single semi-norm \(p_{\iota}\) , and denote by \(E_{\iota}\) the space obtained from E using the topology \(T_{\iota}\) . Then the topology T is the inverse image by the diagonal mapping \(\Delta : E \to \prod_{\iota \in I} E_{\iota}\) of the product topology on \(\prod_{\iota \in I} \mathrm{E}_{\iota}\) (I, p.~9, prop. 7). For each \(\iota \in I\) , write \(\mathbf {N} _ {\iota}\) for the closure of \(\{0 \}\) in \(\mathrm{E} _ {\iota}\) , and by \(F_{\iota} = E_{\iota}/N_{\iota}\) , the normed space defined by the norm \(\hat{p}_{\iota}\) corresponding to \(p_{\iota}\) (II, p.~4, formula (3)); if \(\phi_{\iota}: E_{\iota} \to F_{\iota}\) is the canonical mapping and \(\phi : (x_{\iota}) \mapsto (\phi_{\iota}(x_{\iota}))\) the product mapping, we know that the product topology on \(\prod_{i \in I} E_{i}\) is the inverse image by \(\phi\) of the product topology on \(\prod_{\iota \in I}F_{\iota}\) (GT, II, § 3.9, prop. 18). The topology \(\mathcal{T}\) is, therefore, the inverse image under the composite mapping \(\phi \circ \Delta\) of the product topology on \(\prod_{\mathfrak{t}\in \mathbf{I}}\mathrm{F}_{\mathfrak{t}}\) . In particular, if \(\mathcal{T}\) is Hausdorff then it follows from II, p.~3, prop. 2 that \(\Delta^{-1}(\prod_{\iota \in I} \mathbf{N}_{\iota}) = \{0 \}\) , therefore that the mapping \(\phi \circ \Delta\) is injective, therefore :

PROPOSITION 3. --- Every Hausdorff topological vector space E over K, whose topology is defined by a set of semi-norms, is isomorphic to a sub-space of a product of Banach spaces.

If, further, the topology of E is defined by an enumerable set of semi-norms, then E is metrisable (I, p.~16).

\subsection{Equicontinuity criteria of multilinear mappings for topologies defined by semi-norms}

PROPOSITION 4. --- Let \(E_{i}\) ( \(1 \leqslant i \leqslant n\) ) and F be topological vector spaces over K; we suppose that, for every i, the topology of \(E_{i}\) is defined by a directed set of semi-norms \(\Gamma_{i}\) , and that the topology of F is defined by a set of semi-norms \(\Gamma\) . Then a set H, of multilinear mappings of \(\prod_{i=1}^{n} \mathrm{E}_{i}\) in F is equicontinuous if, and only if, for each semi-norm \(q \in \Gamma\) , and each index \(i\) , there exists a semi-norm \(p_{i} \in \Gamma_{i}\) , and a number \(a > 0\) , such that for each function \(u \in \mathrm{H}\) and point \((x_{i}) \in \prod_{i=1}^{n} \mathrm{E}_{i}\) ,

\[
q \left(u \left(x _ {1}, x _ {2}, \dots , x _ {n}\right)\right) \leqslant a. p _ {1} \left(x _ {1}\right) p _ {2} \left(x _ {2}\right) \dots p _ {n} \left(x _ {n}\right).\tag{5}
\]

The condition is sufficient since it implies that H is equicontinuous at \((0, 0, \ldots, 0)\) and therefore everywhere (I, p.~9, prop. 6).

We show that the condition is necessary. By hypothesis, for every semi-norm \(q \in \Gamma\) and every number \(\beta > 0\) , we have \(q(u(x_1, x_2, ..., x_n)) \leqslant \beta\) for every function \(u \in \mathbf{H}\) provided that \(p_i(x_i) \leqslant \alpha_i\) are true for each index \(i\) , \(1 \leqslant i \leqslant n\) , and certain appropriately chosen numbers \(\alpha_i > 0\) and semi-norms \(p_i \in \Gamma_i\) . As \(\mathbf{K}\) is non-discrete, we can also suppose that, for every \(i\) , we have \(\alpha_i = |\lambda_i| < 1\) where \(\lambda_i \in \mathbf{K}\) . Then let \((x_1, x_2, ..., x_n)\) be any point of \(\prod_{i=1}^{n} E_i\) , and for each index \(i\) , let \(m_i \in \mathbf{Z}\) be an integer such that \(p_i(x_i) \leqslant |\lambda_i|^{m_i + 1}\) ; this can be written as \(p_i(\lambda_i^{-m_i} x_i) \leqslant |\lambda_i| (1 \leqslant i \leqslant n)\) , therefore, by hypothesis, we have

\[
q \left(u \left(x _ {1}, x _ {2}, \dots , x _ {n}\right)\right) \leqslant \beta \left| \lambda_ {1} \right| ^ {m _ {1}} \left| \lambda_ {2} \right| ^ {m _ {2}} \dots \left| \lambda_ {n} \right| ^ {m _ {n}}.\tag{6}
\]

Suppose firstly that one of the \(p_{i}(x_{i})\) is zero, then we can take \(m_{i} \in N\) arbitrarily large, therefore \(q(u(x_{1}, x_{2}, ..., x_{n})) = 0\) . If, on the contrary, all the \(p_{i}(x_{i})\) are \(\neq 0\) , take the integer \(m_{i}\) such that \(|\lambda_{i}|^{m_{i}+2} < p_{i}(x_{i}) \leqslant |\lambda_{i}|^{m_{i}+1}\) for each i; then we have \(|\lambda_{i}|^{m_{i}} < |\lambda_{i}|^{-2}p_{i}(x_{i})\) , from which, by (6), the relation (5) follows with

\[
a = \beta (| \lambda_ {1} | \cdot | \lambda_ {2} | \dots | \lambda_ {n} |) ^ {- 2}. \qquad \text{Q.E.D.}
\]

COROLLARY. --- The set H is equicontinuous if, and only if, for every semi-norm \(q \in \Gamma\) , there exists a neighbourhood of 0 in \(\prod_{i=1}^{n} E_{i}\) , in which the functions \(q \circ u\) , for \(u \in H\) , are uniformly bounded.

The condition is evidently necessary, and the demonstration of prop. 4 shows that it implies an inequality of the form (5) for all \(u \in H\) , and therefore the equicontinuity of H.

We state explicitly the particular case of prop. 4 for linear mappings.

PROPOSITION 5. --- Let E, F be two topological vector spaces over a non-discrete valued division ring K; suppose that the topology of E (resp. F) is defined by a set of semi-norms \(\Gamma\) (resp. \(\Gamma'\) ). Let H be a set of linear mappings of E in F. The following conditions are equivalent :

\begin{enumerate}
\def\labelenumi{\alph{enumi})}
\item
  H is equicontinuous.
\item
  For every semi-norm \(q \in \Gamma'\) , there exists a finite family \((p_i)_{1 \leq i \leq n}\) of semi-norms belonging to \(\Gamma\) and a number \(a > 0\) such that, for all \(x \in E\) and all \(u \in H\) ,
\end{enumerate}

\[
q (u (x)) \leqslant a. \sup _ {1 \leqslant i \leqslant n} p _ {i} (x).\tag{7}
\]

\begin{enumerate}
\def\labelenumi{\alph{enumi})}
\setcounter{enumi}{2}
\tightlist
\item
  For every semi-norm \(q \in \Gamma'\) , the mapping \(\sup_{u \in H} (q \circ u)\) is a continuous semi-norm on \(E\) .
\end{enumerate}

COROLLARY 1. --- Suppose that \(\mathcal{T}, \mathcal{T}'\) are two topologies on a vector space \(\mathbf{E}\) over \(\mathbf{K}\) defined, respectively, by two sets of semi-norms \(\Gamma\) and \(\Gamma'\) . \(\mathcal{T}\) is finer than \(\mathcal{T}'\) if, and only if, for every semi-norm \(q \in \Gamma'\) , there exists a finite family \((p_i)_{1 \leqslant i \leqslant n}\) of semi-norms belonging to \(\Gamma\) and a number \(a > 0\) such that, for all \(x \in \mathbf{E}\) , we have \(q(x) \leqslant a . \sup_{1 \leqslant i \leqslant n} p_i(x)\) .

In fact this shows that the identity mapping of E with topology T, on E with topology \(T'\) , is continuous.

COROLLARY 2. --- Suppose that the topology \(\mathcal{T}\) of a topological vector space \(E\) over \(K\) is defined by a directed set of semi-norms \(\Gamma\) ; for each semi-norm \(p \in \Gamma\) , let \(E_p\) be the space obtained from \(E\) using the topology defined by \(p\) . The set \(E'\) of linear forms on \(E\) that are continuous for \(\mathcal{T}\) is the union of the sets \(E_p'\) , where \(E_p'\) is the set of continuous linear forms in \(E_p\) ( \(p \in \Gamma\) ).

\section{CONVEX SETS}

\subsection{Definition of a convex set}

For any two points x, y of an affine space E, the set of points \(\lambda x + \mu y\) where \(\lambda \geqslant 0\) , \(\mu \geqslant 0\) , \(\lambda + \mu = 1\) is called the closed segment with end points x and y; it reduces to a point when x = y. The complement of x in this segment is called the segment with end points x, y which is open at x and closed at y; it is empty if x = y. Finally the complement of \(\{x, y\}\) in the closed segment with end points x, y is called the open segment with end points x, y; it is empty when x = y.

DEFINITION 1. --- A subset A of an affine space E is convex if, for every two points x, y of A, the closed segment with end points x, y is contained in A.

As \((1 - \lambda)a + \lambda x = a + \lambda(x - a)\) , this definition is equivalent to the following: the set A is convex if, for every point \(a \in A\) , the transform of A by a homothety of centre a and ratio \(\lambda\) where \(0 < \lambda < 1\) , is contained in A (in other words, A is stable for these homotheties).

Examples. --- 1) Every linear affine variety of E (and in particular the empty set) is convex.

\begin{enumerate}
\def\labelenumi{\arabic{enumi})}
\setcounter{enumi}{1}
\item
  The only non-empty convex sets in \(\mathbf{R}\) are the intervals (GT, IV, § 2.4, prop. 1).
\item
  Let E be a vector space and \(\|x\|\) a norm on E; the unit ball B, formed by the points x such that \(\| x\| \leqslant 1\) , is convex since the relations \(\| x\| \leqslant 1\) , \(\| y\| \leqslant 1\) , imply for \(0\leqslant \lambda \leqslant 1\) that
\end{enumerate}

\[
\left\| \lambda x + (1 - \lambda) y \right\| \leqslant \lambda \| x \| + (1 - \lambda) \| y \| \leqslant \lambda + (1 - \lambda) = 1.
\]

Remark. --- Let A be a convex subset of a vector space E; for any scalars \(\alpha > 0\) and \(\beta > 0\) we have \(\alpha A + \beta A = (\alpha + \beta) A\) . In other words, for any \(x \in A\) , \(y \in A\) , there exists \(z \in A\) such that \((\alpha + \beta) z = \alpha x + \beta y\) ; in fact this relation can be written \(z = \frac{\alpha}{\alpha + \beta} x + \frac{\beta}{\alpha + \beta} y\) and we have \(\frac{\alpha}{\alpha + \beta} > 0\) , \(\frac{\beta}{\alpha + \beta} > 0\) and \(\frac{\alpha}{\alpha + \beta} + \frac{\beta}{\alpha + \beta} = 1\) , from which the assertion follows, on using def. 1.

PROPOSITION 1. --- Let \((x_{\iota})\) be a family of points of a convex subset A; every barycentre \(\sum_{\iota} \lambda_{\iota} x_{\iota}\) of the \(x_{\iota}\) formed using positive masses \(\lambda_{\iota}\) (such that \(\sum_{\iota} \lambda_{\iota} = 1\) and \(\lambda_{\iota} = 0\) except for finitely many of the indices, cf.~A, II, §9.3) belongs to A.

Clearly we need only consider the case when the indices are 1, 2, \ldots, p and \(\lambda_{i} > 0\) for each i; the proposition is trivial if p = 1; we prove the result by induction on p.~Put \(\mu = \sum_{i=1}^{p-1} \lambda_{i} > 0\) , and \(y = \sum_{i=1}^{p-1} \frac{\lambda_{i}}{\mu} x_{i}\) ; the induction hypothesis implies that \(y \in A\) . Now as \(\lambda_{p} = 1 - \mu\) and \(\sum_{i=1}^{p} \lambda_{i} x_{i} = \mu y + (1 - \mu) x_{p}\) , its follows from def. 1 that \(\sum_{i=1}^{p} \lambda_{i} x_{i}\) belongs to A.

PROPOSITION 2. --- Let E and F be two affine spaces and f be an affine linear mapping of E in F; then the image of a convex subset of E under f, and the inverse image of a convex subset of F under f are both convex.

The image under f of the closed segment with end points x, y is the closed segment with end points \(f(x)\) , \(f(y)\) , hence the first statement. We deduce that the inverse image of a closed segment of F under f contains each closed segment whose end points belong to it; the second statement of prop. 2 follows.

In particular the image of a convex set under a homothety or a translation is a convex set.

PROPOSITION 3. --- In the affine space E, let H be a hyperplane defined by the relation \(g(x) = 0\) , where g is a non-constant affine function on E. Then the half-spaces defined by the relations \(g(x) \geqslant 0\) , \(g(x) \leqslant 0\) , \(g(x) > 0\) , \(g(x) < 0\) are convex.

For these are the inverse images under g of intervals of R and thus are convex.

With the notations of prop. 3 the points of a subset M of an affine space are on the same side (resp. strictly on the same side) of the hyperplane H if M is contained in one of the half-spaces defined by \(g(x) \geqslant 0\) , \(g(x) \leqslant 0\) (resp. \(g(x) > 0\) or \(g(x) < 0\) ).

PROPOSITION 4. --- The points of A, a convex subset of an affine space E are strictly on the same side of a hyperplane H if, and only if, A does not meet H.

Clearly the condition is necessary. Conversely suppose that it is satisfied and let \(g(x) = 0\) , be an equation defining H (g is an affine linear mapping of E in R). The set \(g(\mathbf{A})\) is convex in R, therefore it is an interval, and \(0 \notin g(\mathbf{A})\) . Hence \(g(x)\) is of fixed sign for all \(x \in A\) .

\subsection{Intersections of convex sets. Products of convex sets}

PROPOSITION 5. --- The intersection of any family of convex subsets of an affine space E is convex.

The proposition follows immediately from def. 1 of II, p.~7.

PROPOSITION 6. --- Let \((\mathrm{E}_{\iota})_{\iota\in\mathbf{I}}\) be a family of vector spaces, and for each \(\iota\in I\) , let \(A_{\iota}\) be a non-empty subset of \(E_{\iota}\) . Then the set \(A=\prod_{\iota\in I}A_{\iota}\) is convex in \(E=\prod_{\iota\in I}E_{\iota}\) , if, and only if, for all \(\iota\in I\) , the set \(A_{\iota}\) is convex in \(E_{\iota}\) .

In fact, each projection \(\mathrm{pr}_{\iota}\) is a linear mapping and we have \(A_{\iota} = \mathrm{pr}_{\iota}A\) and \(A = \bigcap_{\iota \in I} \mathrm{pr}_{\iota}^{-1}(A_{\iota})\) ; the proposition follows from props. 2 and 5 above.

COROLLARY. --- In the space \(\mathbf{R}^n\) every parallelotope (GT, VI, § 1.3) is a convex subset.

For it is the image under an affine linear mapping of a rectangular parallelepiped, and this last is convex by prop. 6.

PROPOSITION 7. --- Let A and B be two convex subsets of the vector space E. For any real numbers \(\alpha\) , \(\beta\) the set \(\alpha A + \beta B\) (set of points of the form \(\alpha x + \beta y\) , where x varies in A, and y in B) is convex.

For \(\alpha A + \beta B\) is the image of the convex subset \(A \times B\) of \(E \times E\) under the linear mapping \((x, y) \mapsto \alpha x + \beta y\) of \(E \times E\) in E.

\subsection{Convex envelope of a set}

DEFINITION 2. --- Given a subset A of an affine space E, we call the intersection of all convex sets containing A, the convex envelope of A, that is to say (II, p.~9, prop. 5) it is the smallest convex set containing A.

PROPOSITION 8. --- For any family \((\mathbf{A}_{\iota})_{\iota \in \mathbf{I}}\) of convex subsets of an affine space E, the convex envelope of \(\bigcup_{\iota \in \mathbf{I}}\mathbf{A}_{\iota}\) is precisely the set of linear combinations \(\sum_{\iota \in \mathbf{I}}\lambda_{\iota}x_{\iota}\) , where \(x_{\iota}\in \mathbf{A}_{\iota},\lambda_{\iota}\geqslant 0\) for all \(\iota \in \mathrm{I}\) ( \(\lambda_{\iota} = 0\) except for finitely many indices) and \(\sum_{\iota \in \mathbf{I}}\lambda_{\iota} = 1\) .

Denote the set of these linear combinations by C, clearly C is contained in every convex set which contains all the \(A_{\iota}\) (II, p.~8, prop. 1); on the other hand \(A_{\iota} \subset C\) for every \(\iota\) . All that remains to be proved is that C is convex. Let \(x = \sum_{i} \lambda_{i} x_{i}\) , \(y = \sum_{i} \mu_{i} y_{i}\) be two points of C and \(\alpha\) be a number such that \(0 < \alpha < 1\) , write \(\gamma_{\iota} = \alpha \lambda_{\iota} + (1 - \alpha)\mu_{\iota}\) for every \(\iota \in \mathbf{I}\) , and let J be the set (finite) of the indices of I for which \(\gamma_{\iota} \neq 0\) . We can write \(\alpha x + (1 - \alpha)y = \sum_{\iota \in \mathbf{J}}\gamma_{\iota}z_{\iota}\) , where

\[
z _ {\iota} = \gamma_ {\iota} ^ {- 1} (\alpha \lambda_ {\iota} x _ {\iota} + (1 - \alpha) \mu_ {\iota} y _ {\iota})
\]

belongs to \(\mathbf{A}_{\iota}\) for all \(\iota \in \mathbf{J}\) ; but \(\sum_{\iota \in \mathbf{J}} \gamma_{\iota} = \alpha \sum_{\iota \in \mathbf{I}} \lambda_{\iota} + (1 - \alpha) \sum_{\iota \in \mathbf{I}} \mu_{\iota} = 1\) , and we see that \(\alpha x + (1 - \alpha)y \in \mathbf{C}\) . The proposition is proved.

COROLLARY 1. --- The convex envelope of a subset A of E is identical with the set of linear combinations \(\sum_{i} \lambda_{i} x_{i}\) , where \((x_{i})\) is any finite family of points of A, the numbers \(\lambda_{i} > 0\) , for all \(i\) and \(\sum_{i} \lambda_{i} = 1\) .

The dimension of the affine linear variety (A, II, § 9.3) generated by the convex set A is called the dimension of A.

Let E be a vector space. The convex envelope C, of the balanced envelope of a set A in E is called the balanced convex envelope (or the symmetric convex envelope) of A; clearly it is the smallest symmetric convex set that contains A; it is also the convex envelope of \(A \cup (-A)\) , since every point of the balanced envelope of A belongs to a segment with extremities a and -a where \(a \in A\) . The set C coincides with the set of linear combinations \(\sum_{i} \lambda_{i} x_{i}\) where \(x_{i} \in A\) and \(\sum_{i} |\lambda_{i}| \leqslant 1\) ; for it is clear

that this set of points is convex and contains A and -A; it is sufficient to prove that it is contained in C, and for this we need consider only those linear combinations for which \(\mu = \sum_{i} |\lambda_{i}| > 0\) ; we can then write \(\sum_{i} \lambda_{i} x_{i} = \mu \cdot \sum_{i} \alpha_{i} y_{i}\) with \(\alpha_{i} = \lambda_{i} / \mu\) and \(y_{i} = x_{i}\) , if \(\lambda_{i} \geqslant 0\) ; and \(\alpha_{i} = -\lambda_{i} / \mu\) ; \(y_{i} = -x_{i}\) if \(\lambda_{i} < 0\) ; clearly \(\sum_{i} \alpha_{i} = 1\) , and our assertion is proved.

COROLLARY 2. --- Let f be an affine linear mapping of the affine space E in the affine space F; for each subset A of E, the convex envelope of \(f(A)\) is the image under f of the convex envelope of A.

There is a similar statement for linear mappings and balanced convex envelopes.

\subsection{Convex cones}

DEFINITION 3. --- A subset C of an affine space E is a cone with vertex \(x_0\) if C is invariant for all homotheties of centre \(x_0\) and ratio \(> 0\) .

We shall suppose in this No.~and in the one following, that we have chosen the vertex of the cone being considered, as origin in E; i.e.~we suppose that E is a vector space, and when we speak of a cone, it is to be understood that this cone has vertex 0. The set of points of the form \(\lambda a\) for \(\lambda > 0\) (resp. \(\lambda \geqslant 0\) ), where \(a\) is a non-null vector, is called an open half line (resp. closed half-line) originating at 0.

A cone C of vertex 0 is said to be pointed if \(0 \in C\) , and non-pointed otherwise. A pointed cone is either the single point \(\{0\}\) or is the union of a set of closed half-lines originating at 0. A non-pointed cone is the union (possibly empty) of open half lines originating at 0. If C is a non-pointed cone, then \(C \cup \{0\}\) is a pointed cone. If C is a pointed cone, then \(C - \{0\}\) is a non-pointed cone.

If C is a non-pointed convex cone, then \(C \cup \{0\}\) is a pointed convex cone. However, if C is a pointed convex cone, \(C - \{0\}\) is not necessarily convex. We say that a pointed convex cone is proper if it does not contain any line passing through 0. Then

PROPOSITION 9. --- A pointed convex cone C is proper if and only if the non-pointed cone C', which is the complement of 0 in C, is convex.

If C contains a line through 0 then clearly \(C'\) is not convex. Suppose now that C is proper and let x, y be two points of \(C'\) . The closed segment with end points x, y is contained in C; if it contains 0 then \(\lambda x + (1 - \lambda)y = 0\) for some \(\lambda\) with \(0 < \lambda < 1\) , therefore \(x = \mu y\) with \(\mu < 0\) . Thus C contains the line through 0 and x, contrary to hypothesis.

PROPOSITION 10. --- A subset C of E is a convex cone if and only if \(C + C \subset C\) and \(\lambda C \subset C\) for all \(\lambda > 0\) .

For the condition \(\lambda C \subset C\) for all \(\lambda > 0\) characterises the cones. If C is convex we have \(C + C = \frac{1}{2}C + \frac{1}{2}C = C\) (II, p.~8, Remark). Conversely, if the cone C is such that \(C + C \subset C\) , then for \(0 < \lambda < 1\) , we have \(\lambda C + (1 - \lambda)C = C + C \subset C\) , which shows that C is convex.

COROLLARY 1. --- If C is a non-empty convex cone, the vector space generated by C is the set C - C (the set of points x - y where x, y vary in C).

For, if V = C - C, then V is non empty, we have \(\lambda V = V\) for all \(\lambda \neq 0\) , and \(V + V = C + C - (C + C) \subset C - C = V\) , which shows that V is a vector subspace. Finally every vector subspace that contains C also contains V.

COROLLARY 2. --- If C is a pointed convex cone, the largest vector subspace contained in C is the set \(C \cap (-C)\) .

For, if \(W = C \cap (-C)\) , then W is non-empty and \(\lambda W = W\) for all \(\lambda \neq 0\) , also

\[
\mathrm{W} + \mathrm{W} \subset (\mathrm{C} + \mathrm{C}) \cap (- (\mathrm{C} + \mathrm{C})) \subset \mathrm{C} \cap (- \mathrm{C}) = \mathrm{W},
\]

which shows that W is a vector subspace. Clearly every vector subspace contained in C is also contained in W.

Obviously, if f is a linear mapping of E in a vector space F, then \(f(\mathbf{C})\) , the image of a convex cone C in E, is a convex cone in F. Every intersection of convex cones (with vertex 0) in E is a convex cone. For every subset A of E the intersection of convex cones containing A (these exist, E itself is one such cone) is the smallest convex cone that contains A; it is called the convex cone generated by A.

PROPOSITION 11. --- Let \((\mathbf{C}_{\iota})_{\iota \in \mathbf{I}}\) be a family of convex cones in E; the convex cone generated by the union of the \(\mathbf{C}_{\iota}\) is identical with the set of points \(\sum_{\iota \in \mathbf{J}} x_{\iota}\) , where \(\mathbf{J}\) is any finite

subset of I and \(x_{\iota} \in \mathbf{C}_{\iota}\) for all \(\iota \in \mathbf{J}\) .

In fact, it is obvious that C, the set of such points, is a convex cone containing the union of the \(C_{\iota}\) , and that it is contained in any convex cone which contains this union.

COROLLARY. --- For any subset A of E, the convex cone generated by A, is identical with the set of linear combinations \(\sum_{i\in J}\lambda_{i}x_{i}\) , where \((x_{i})_{i\in J}\) is any finite non-empty family of points of A, and where \(\lambda_{i}>0\) for all \(i\in J\) .

It is sufficient to see that, if a convex cone contains a point \(x \neq 0\) of A then it also contains the half-line \(C_{x}\) of the points \(\lambda x\) where \(\lambda\) varies in the set of positive numbers and that \(C_{x}\) is a convex cone.

PROPOSITION 12. --- If A is a convex set in E, then the convex cone generated by A is identical with \(C = \bigcup_{\lambda > 0} \lambda A\) .

The set C is clearly a cone; it is sufficient to show that C is convex. Let \(\lambda x, \mu y\) be two points of C ( \(\lambda > 0, \mu > 0, x \in A, y \in A\) ). Let \(\alpha, \beta\) be two numbers \textgreater{} 0 such that \(\alpha + \beta = 1\) . Then \(\alpha \lambda x + \beta \mu y = (\alpha \lambda + \beta \mu) z\) , with \(z \in A\) , and \(\alpha \lambda + \beta \mu > 0\) ; hence \(\alpha \lambda x + \beta \mu y \in C\) .

Remarks. --- 1) With the hypotheses of prop. 12, if \(0 \notin \mathbf{A}\) , then the cone \(C\) is non-pointed, thus \(C \cup \{0\}\) is proper.

\begin{enumerate}
\def\labelenumi{\arabic{enumi})}
\setcounter{enumi}{1}
\tightlist
\item
  Let A be any convex set in E; consider the convex set \(A_{1} = A \times \{1\}\) in the space \(F = E \times R\) and the convex cone C with vertex 0 that is generated by \(A_{1}\) . Prop. 12 shows that \(A_{1}\) is the intersection of C and of the hyperplane \(E \times \{1\}\) in F. Every convex set in E can, therefore, be considered as the projection on E of the intersection of a convex cone with vertex 0 in F and the hyperplane \(E \times \{1\}\) .
\end{enumerate}

\subsection{Ordered vector spaces}

A preorder structure, on a vector space E, denoted by \(x \preccurlyeq y\) or \(y \succcurlyeq x\) , is compatible with the vector space structure of E if it satisfies the following two axioms;

\((\mathrm{EO}_{\mathrm{I}})\) If \(x \preccurlyeq y\) then \(x + z \preccurlyeq y + z\) for all \(z \in \mathrm{E}\) .

\((\mathrm{EO}_{\mathrm{II}})\) If \(x \succcurlyeq 0\) then \(\lambda x \succcurlyeq 0\) for every scalar \(\lambda \geqslant 0\) .

The vector space E, carrying these two structures, is called a preordered vector space (resp. an ordered vector space when the relation of preorder on E is an order).

Note that axiom \((EO_{I})\) means that the preorder structure and the additive group structure of E are compatible, that is to say, E carrying these two structures, is a preordered group (A, VI, p.~3).

Example. --- On the space \(E = R^{A}\) of all finite real-valued functions defined over A, the relation of order given by « for all \(t \in A\) , \(x(t) \leqslant y(t)\) » is compatible with the vector space structure of E.

PROPOSITION 13. --- (i) The set P, of elements \(\succcurlyeq 0\) , of a preordered vector space E, is a pointed convex cone.

\begin{enumerate}
\def\labelenumi{(\roman{enumi})}
\setcounter{enumi}{1}
\item
  Conversely, if \(\mathbf{P}\) is a pointed convex cone in \(\mathbf{E}\) , then the relation \(y - x \in \mathbf{P}\) is a preorder relation on \(\mathbf{E}\) , and the preorder structure that it defines is the only one that is compatible with the vector space structure of E and for which P is the set of elements \(\succcurlyeq 0\) .
\item
  The relation \(y - x \in \mathbf{P}\) , with \(\mathbf{P}\) a pointed convex cone, is an order relation on \(\mathbf{E}\) if and only if \(\mathbf{P}\) is a proper cone.
\item
  Axioms \((\mathrm{EO}_{\mathrm{I}})\) and \((\mathrm{EO}_{\mathrm{II}})\) imply \(P + P \subset P\) and \(\lambda P \subset P\) for all \(\lambda > 0\) . As \(0 \in P\) , it follows that P is a pointed convex cone (II, p.~11, prop. 10).
\item
  Conversely, if \(\mathbf{P}\) is a pointed convex cone, the relation \(\mathbf{P} + \mathbf{P} \subset \mathbf{P}\) implies that the relation \(y - x \in \mathbf{P}\) is a preorder compatible with the additive group structure of E (A, VI, p.~3, prop. 3); clearly writing it \(x \preccurlyeq y\) , the set \(\mathbf{P}\) is identical with the set of \(x \geqslant 0\) ; further the relation \(\lambda \mathbf{P} \subset \mathbf{P}\) for all \(\lambda \geqslant 0\) shows that axiom \((\mathrm{EO}_{\mathrm{II}})\) is satisfied. (iii) To say that \(\mathbf{P}\) is proper means that \(\mathbf{P} \cap (-\mathbf{P}) = \{0\}\) (II, p.~11, cor. 2), hence that \(y - x \in \mathbf{P}\) is an order relation.
\end{enumerate}

Example. --- * Let H be a real Hilbert space; in the vector space \(\mathcal{L}(\mathrm{H})\) of continuous endomorphisms of H, the positive hermitian endomorphisms form a proper pointed convex cone; this cone, therefore, defines an order structure compatible with the vector space structure of \(\mathcal{L}(\mathrm{H})\) and for which the relation \(A \leqslant B\) means that \(B - A\) is a positive hermitian endomorphism. *

For any pointed convex cone P in the vector space E, the set \(P \cap (-P)\) is a vector subspace, H, of E (II, p.~11, cor. 2). The canonical image \(P'\) of P in E/H is a convex cone and the inverse image of \(P'\) in E is P. Thus \(P' \cap (-P') = \{0\}\) , and \(P'\) defines an order structure on E/H that is compatible with its vector space structure.

A linear form f on a preordered vector space E is said to be positive if \(x \geqslant 0\) in E implies \(f(x) \geqslant 0\) . Or, alternatively, if the convex cone P of elements \(\geqslant 0\) in E is contained in the half space of those x for which \(f(x) \geqslant 0\) . Clearly, in the dual E* to E, the set of positive linear forms is a pointed convex cone.

\subsection{Convex cones in topological vector spaces}

PROPOSITION 14. --- In a topological vector space E, the closure of a convex set (resp. of a convex cone) is a convex set (resp. a convex cone with the same vertex).

For, let A be a convex set; the mapping \((x, y) \mapsto \lambda x + (1 - \lambda)y\) , where \(0 < \lambda < 1\) , is continuous in E × E and maps A × A in A; thus (GT, I, §2.1, th. 1) it maps \(\overline{A} \times \overline{A}\) in \(\overline{A}\) , which shows that \(\overline{A}\) is convex. Similarly, if C is a convex cone with vertex 0 then \(\overline{C} + \overline{C} \subset \overline{C}\) and \(\lambda\overline{C} \subset \overline{C}\) for all \(\lambda > 0\) .

DEFINITION 4. --- For any set A of a topological vector space E, the intersection of all the closed convex sets containing A is called the convex closed envelope of A; it is the smallest convex closed set containing A.

From prop. 14, the convex closed envelope of A is the closure of the convex envelope of A; it is clearly the same as the convex closed envelope of \(\overline{A}\) .

Similarly we call the smallest symmetric, convex, closed set that contains A, the symmetric convex closed envelope (or the balanced convex closed envelope) of A; it is the closure of the symmetric convex envelope of A (II, p.~10); it is also the symmetric convex closed envelope of \(\overline{A}\) .

PROPOSITION 15. --- Let \(A_{i}\) ( \(1 \leqslant i \leqslant n\) ) be a finite number of compact convex sets in a Hausdorff topological vector space E. Then the convex envelope of the union of the \(A_{i}\) is compact (and is, therefore, the same as the convex closed envelope of this union).

Let B be the compact set in \(R^{n}\) defined by the points \((\lambda_{1},\lambda_{2},\ldots,\lambda_{n})\) where \(\lambda_{i}\geqslant0(1\leqslant i\leqslant n)\) , and \(\sum_{i=1}^{n}\lambda_{i}=1\) . Define a continuous mapping of \(B\times\prod_{i=1}^{n}A_{i}\subset R^{n}\times E^{n}\) in E by the formula \((\lambda_{1},\lambda_{2},\ldots,\lambda_{n},x_{1},x_{2},\ldots,x_{n})\mapsto\sum_{i=1}^{n}\lambda_{i}x_{i}\) . The convex envelope C of \(\bigcup_{i=1}^{n}A_{i}\) is the image of \(B\times\prod_{i=1}^{n}A_{i}\) under this mapping; as \(B\times\prod_{i=1}^{n}A_{i}\) is compact and E is Hausdorff, it follows that C is compact.

COROLLARY 1. --- In a Hausdorff topological vector space the convex envelope of a finite set is compact.

COROLLARY 2. --- In a topological vector space E, the convex envelope of a finite set is precompact.

In fact, let \(j\) be the canonical mapping of \(\mathbf{E}\) in its Hausdorff completion \(\hat{\mathbf{E}}\) ; if \(\mathbf{C}\) is the convex envelope of \(\mathbf{A}\) , then \(j(\mathbf{C})\) is the convex envelope of the finite set \(j(\mathbf{A})\) in \(\hat{\mathbf{E}}\) , hence \(j(\mathbf{C})\) is compact (cor. 1) and therefore \(\mathbf{C}\) is precompact (GT, II, § 4.2).

PROPOSITION 16. --- Let A be a convex subset, with at least one interior point \(x_{0}\) , of a topological vector space E. For any point \(x \in \overline{A}\) , every point of the open segment with end points \(x_{0}\) , x lies in the interior of A.

For any point y of this segment, let f be the homothety of centre y and ratio \(\lambda < 0\) , which transforms \(x_{0}\) into x. If V is an open neighbourhood of \(x_{0}\) contained in A, then \(f(\mathrm{V})\) is a neighbourhood of x and therefore contains a point \(f(z) \in \mathrm{A}\) ; now

\[
f (z) - y = \lambda (z - y) = \lambda (z - f (z)) + \lambda (f (z) - y),
\]

hence \(y - f(z) = \frac{\lambda}{\lambda - 1} (z - f(z))\) , so that \(y\) is transformed into \(z\) by the homothety \(g\) , of centre \(f(z)\) and ratio \(\mu = \lambda / (\lambda - 1)\) ; since \(0 < \mu < 1\) , \(g\) transforms \(V\) into a neighbourhood of 0 contained in \(A\) . The proposition is proved.

COROLLARY 1. --- The interior \(\mathring{\mathrm{A}}\) of a convex set \(\mathrm{A}\) , is itself a convex set; if \(\mathring{\mathrm{A}}\) is not empty, then it coincides with the interior of \(\overline{\mathrm{A}}\) , and \(\overline{\mathrm{A}}\) is a convex set that coincides with the closure of \(\mathring{\mathrm{A}}\) .

It follows from prop. 16, that if \(\mathring{\mathrm{A}}\) is not empty, then it is a convex set and every point of \(\overline{\mathrm{A}}\) is a cluster point of \(\mathring{\mathrm{A}}\) . Next we show that every interior point of \(\overline{\mathrm{A}}\) belongs to \(\mathring{\mathrm{A}}\) . Let \(x\) be an interior point of \(\overline{\mathrm{A}}\) and suppose, for definiteness that \(x = 0\) . Let \(\mathrm{V}\) be a symmetric neighbourhood of 0 that is contained in \(\overline{\mathrm{A}}\) and let \(y \in \mathring{\mathrm{A}} \cap \mathrm{V}\) ; now \(-y \in \overline{\mathrm{A}}\) , and therefore, by prop. 16, we see that \(0 \in \mathring{\mathrm{A}}\) , if \(y \neq 0\) ; this is obviously true if \(y = 0\) .

COROLLARY 2. --- The interior \(\mathring{\mathrm{C}}\) of a convex cone \(\mathrm{C}\) , is itself a convex cone; if \(\mathring{\mathrm{C}}\) is not empty then it coincides with the interior of \(\overline{\mathrm{C}}\) , and \(\overline{\mathrm{C}}\) is a pointed convex cone that coincides with the closure of \(\mathring{\mathrm{C}}\) .

Since homotheties of ratio \textgreater{} 0 and centre 0 transform C into itself, they do the same for \(\overset{\circ}{C}\) , thus \(\overset{\circ}{C}\) is a cone; the remainder of the corollary follows from cor. 1 and the obvious remark that if C is not empty then \(\overline{C}\) contains the vertex of C.

Let H be a closed hyperplane in the topological vector space E over R; it has an equation of the form \(f(x) = \alpha\) , where f is a continuous linear form that is not identically zero in E (I, p.~13, th. 1). The closed half spaces defined respectively by \(f(x) \leqslant \alpha\) and \(f(x) \geqslant \alpha\) are therefore closed convex sets; their complements defined respectively by \(f(x) > \alpha\) and \(f(x) < \alpha\) , are open convex sets. We say that these half-spaces are the closed (resp. open) half spaces determined by H.

PROPOSITION 17. --- In a topological vector space E, let A be a set with at least one interior point, and such that all its points lie on the same side of an hyperplane H. Then H is closed, the interior points of A lie strictly on the same side of H, and the cluster points of A lie on the same side of H. In particular open (resp. closed) half spaces are determined by closed hyperplanes.

In fact suppose that H contains the origin and that \(f(x) = 0\) is an equation of H; suppose, for definiteness, that \(f(x) \geqslant 0\) for all \(x \in A\) . The half space formed by the points y such that \(f(y) > -1\) contains at least one interior point, and, by translation, the same is true of the half space of points y such that \(f(y) > 0\) ; this shows that H is closed (I, p.~11, corollary). Then we know that f is a strict morphism of E on R (I, p.~13, corollary), therefore \(f(\mathring{\mathrm{A}})\) is an open set in R. This set cannot contain 0 or it would contain numbers \textless{} 0 contrary to hypothesis; it is thus contained in the open interval \(\bigg]0, +\infty[\) . On the other hand, the half space of those y for which \(f(y) \geqslant 0\) is closed and contains A, therefore it contains \(\overline{A}\) .

COROLLARY. --- Let P be a pointed convex cone, with at least one interior point, of the topological vector space E. Then each linear form f that is not identically zero on E, and is positive for the preorder structure defined by P(II, p.~13), is necessarily continuous. Further, if x is interior to P then \(f(x) > 0\) and if x is a cluster point of P then \(f(x) \geqslant 0\) .

Apply prop. 17 to the case \(\mathbf{A} = \mathbf{P}\) where \(\mathbf{H}\) is the hyperplane with the equation \(f(x) = 0\) .

Remark. --- In a topological vector space E, every convex set C is connected. In fact, if \(a \in C\) , then C is a union of segments with end point a and closed at a; these are connected and the result follows from GT, I, § 11.1, prop. 2.

\subsection{Topologies on ordered vector spaces}

Let E be an ordered vector space. A topology on E is compatible with the ordered vector space structure of E if it is both compatible with the vector space structure of E and subject to the following axiom :

(TO) The convex cone of the \(x\) with \(x \geqslant 0\) , is closed in E.

An ordered vector space E with a compatible topology is called an ordered topological vector space.

Examples. --- The space \(R^{n}\) with its usual topology and the order structure that is the product of the order structure of its factors is an ordered topological vector space. On the other hand, for \(n \geqslant 2\) , when \(R^{n}\) carries the lexicographical order (S, III, § 2.6), the usual topology is not compatible with the ordered vector space structure of \(R^{n}\) .

Let \(\mathbf{A}\) be a set; the vector space \(\mathcal{B}(\mathbf{A};\mathbf{R})\) of real valued bounded functions defined on \(\mathbf{A}\) , with the topology defined by the norm \(\| x\| = \sup_{t\in \mathbf{A}}|x(t)|\) and the order structure

induced by the product order structure of \(R^{A}\) , is an ordered topological vector space.

In an ordered topological vector space E, the set of elements \(x \leqslant 0\) is closed; since translations are homeomorphisms, we deduce that, for all \(a \in E\) , the set of elements \(x \geqslant a\) (resp. \(x \leqslant a\) ) is closed. Since \(\{0\}\) is the intersection of the sets \(x \geqslant 0\) and \(x \leqslant 0\) , it follows that \(\{0\}\) is closed and that E is Hausdorff.

PROPOSITION 18. --- In an ordered topological vector space E, let H be a set directed by the relation \(\leqslant\) . If the section filter of H has a limit in E, then this limit is the upper bound of H.

For, let \(b = \lim_{x\in \mathbf{H}}x\) ; for every \(y\in \mathbf{H}\) , the set of \(x\in \mathbf{H}\) such that \(x\geqslant y\) is a set of

the section filter of H, therefore b is a cluster point of this set; but as the set \(x \geqslant y\) is closed in E, we have \(b \geqslant y\) , thus b is an upper bound of H. On the other hand, if a is an upper bound of H, then H is contained in the closed set \(x \leqslant a\) ; as b is a cluster point of H, we have \(b \leqslant a\) , which completes the proof (II, p.~72, exerc. 42).

\subsection{Convex functions}

DEFINITION 5. --- Let X be a convex subset of the affine space E. A real-valued finite function, defined over X is convex (resp. strictly convex) if for any two distinct points x, y of X and any real number \(\lambda\) , \(0 < \lambda < 1\) , we have :

\[
f (\lambda x + (1 - \lambda) y) \leqslant \lambda f (x) + (1 - \lambda) f (y)\tag{1}
\]

(resp.

\[
f (\lambda x + (1 - \lambda) y) <   \lambda f (x) + (1 - \lambda) f (y)).\tag{2}
\]

When E = R, this definition of convex function is the same as that in FVR, I, p.~32. Further, f is convex (resp. strictly convex) in X if, and only if, for every affine line \(D \subset E\), the restriction of f to \(X \cap D\) is convex (resp. strictly convex) in \(X \cap D\).

Examples. --- If \(f\) is an affine linear function on \(\mathbf{E}\) , then \(f\) and \(f^2\) are convex functions on \(\mathbf{E}\) ; this is obvious for \(f\) since

\[
f (\lambda x + (1 - \lambda) y) = \lambda f (x) + (1 - \lambda) f (y);
\]

on the other hand, if \(\alpha = f(x)\) , \(\beta = f(y)\) , then;

\[
\lambda \alpha^ {2} + (1 - \lambda) \beta^ {2} - (\lambda \alpha + (1 - \lambda) \beta) ^ {2} = \lambda (1 - \lambda) (\alpha - \beta) ^ {2} \geqslant 0
\]

for \(0 < \lambda < 1\) ; further, the restriction of \(f^2\) to an affine line \(\mathbf{D} \subset \mathbf{E}\) is strictly convex if \(f|\mathbf{D}\) is not a constant.

A real-valued function \(f\) , defined over \(\mathbf{X}\) , is concave (resp. strictly concave) if \(-f\) is convex (resp. strictly convex). That is to say, for every two distinct points \(x, y\) of \(\mathbf{X}\) and every number \(\lambda\) , such that \(0 < \lambda < 1\) , we have

\[
f (\lambda x + (1 - \lambda) y) \geqslant \lambda f (x) + (1 - \lambda) f (y)
\]

(resp.

\[
f (\lambda x + (1 - \lambda) y) > \lambda f (x) + (1 - \lambda) f (y)).
\]

A mapping of X in R is affine if it is both convex and concave (cf.~II, p.~78, exerc. 11).

PROPOSITION 19. --- Let X be a convex set of the affine space E; and let f be a real-valued function defined over X. Denote the set of points \((x, a) \in \mathbf{E} \times \mathbf{R}\) for which \(x \in X\) and \(f(x) \leqslant a\) (resp. \(x \in X\) and \(f(x) < a\) ) by F (resp. F'). Then the following conditions are equivalent:

\begin{enumerate}
\def\labelenumi{\alph{enumi})}
\item
  The function \(f\) is convex.
\item
  The set \(\mathbf{F}\) in the affine space \(\mathbf{E} \times \mathbf{R}\) is convex.
\item
  The set \(\mathbf{F}'\) in the affine space \(\mathbf{E} \times \mathbf{R}\) is convex.
\end{enumerate}

We show that \(a) \Rightarrow c)\) . Let \((x, a)\) and \((y, b)\) be two points of \(F'\) and \(0 < \lambda < 1\) , then \(f(x) < a\) , \(f(y) < b\) and if \(f\) is convex

\[
f (\lambda x + (1 - \lambda) y) \leqslant \lambda f (x) + (1 - \lambda) f (y) <   \lambda a + (1 - \lambda) b
\]

which shows that the point \(\lambda(x, a) + (1 - \lambda)(y, b)\) of \(\mathbf{E} \times \mathbf{R}\) belongs to \(\mathbf{F}'\) . Thus \(\mathbf{F}'\) is convex.

Next we show that \(c) \Rightarrow b)\) . If \((x, a), (y, b)\) are two points of \(F\) then for every \(\varepsilon > 0\) , \((x, a + \varepsilon)\) and \((y, b + \varepsilon)\) belong to \(F'\) and, if \(0 < \lambda < 1\) , the same is true of \((\lambda x + (1 - \lambda)y, \lambda a + (1 - \lambda)b + \varepsilon)\) ; by the definition of \(F\) this implies that \((\lambda x + (1 - \lambda)y, \lambda a + (1 - \lambda)b)\) belongs to \(F\) .

Finally \(b) \Rightarrow a)\) , for (with the above notation), if \((\lambda x + (1 - \lambda)y, \lambda a + (1 - \lambda)b)\) belongs to F then

\[
f (\lambda x + (1 - \lambda) y) \leqslant \lambda a + (1 - \lambda) b
\]

provided \(a \geqslant f(x)\) and \(b \geqslant f(y)\) ; hence (1) follows and \(f\) is convex.

COROLLARY. --- If \(f\) is convex in \(\mathbf{X}\) , then for all \(\alpha \in \mathbf{R}\) , the set of \(x \in \mathbf{X}\) such that \(f(x) \leqslant \alpha\) (resp. \(f(x) < \alpha\) ) is convex.

In fact, it is the projection on E of the intersection of F (resp. F') and the hyperplane \(E \times \{\alpha\}\) in \(E \times R\) .

PROPOSITION 20. --- Let f be a convex function, defined over a convex set X of the affine space E. Then for every family \((x_{i})_{1\leqslant i\leqslant p}\) of \(p\) points of X and every family \((\lambda_{i})_{1\leqslant i\leqslant p}\) . of \(p\) real numbers, all \(\geqslant 0\) , such that \(\sum_{i = 1}^{p}\lambda_{i} = 1\) , we have:

\[
f \left(\sum_ {i = 1} ^ {p} \lambda_ {i} x _ {i}\right) \leqslant \sum_ {i = 1} ^ {p} \lambda_ {i} f \left(x _ {i}\right).\tag{3}
\]

If \(f\) is strictly convex and if \(\lambda_i > 0\) for all \(i\) , then

\[
f \left(\sum_ {i = 1} ^ {p} \lambda_ {i} x _ {i}\right) <   \sum_ {i = 1} ^ {p} \lambda_ {i} f \left(x _ {i}\right),\tag{4}
\]

unless all the \(x_{i}\) are equal.

The inequality (3) follows from II, prop. 19 above and II, p.~8, prop. 1. Suppose that the \(x_{i}\) are not all equal (which implies \(p \geqslant 2\) ) and that the \(\lambda_{i}\) are all \(>0\) ; then the point \(z = \sum_{i=1}^{p} \lambda_{i} x_{i}\) differs from at least one \(x_{i}\) . Suppose for definiteness that \(z \neq x_{1}\) , write \(z = \lambda_{1} x_{1} + (1 - \lambda_{1}) y_{1}\) where \(y_{1} = \sum_{i=2}^{p} \frac{\lambda_{i}}{1 - \lambda_{1}} x_{i}\) . Then \(y_{1} \neq x_{1}\) and, as \(0 < \lambda_{1} < 1\) , we have, by hypothesis,

\[
f (z) <   \lambda_ {1} f (x _ {1}) + (1 - \lambda_ {1}) f (y _ {1}).
\]

But by (3) \(f(y_{1}) \leqslant \sum_{i=2}^{p} \frac{\lambda_{i}}{1 - \lambda_{1}} f(x_{i})\) , and the inequality (4) follows.

\subsection{Operations on convex functions}

Let \(\mathbf{X}\) be a convex set of an affine space \(\mathbf{E}\) . If \(f_{i}(1 \leqslant i \leqslant p)\) are finitely many convex functions defined over \(\mathbf{X}\) and \(c_{i} (1 \leqslant i \leqslant p)\) are numbers \(\geqslant 0\) then the function \(f = \sum_{i=1}^{p} c_{i} f_{i}\) is convex over \(\mathbf{X}\) .

If \((f_{i})\) is any family of convex functions defined over \(\mathbf{X}\) and if \(g\) , the upper envelope of the family in \(\mathbf{X}\) , is finite then \(g\) is convex.

Finally if H is a set of convex functions defined over X, and F is a filter on H that converges simply in X to the finite real valued function \(f_{0}\) , then \(f_{0}\) is convex over X.

\subsection{Convex functions over an open convex set}

PROPOSITION 21. --- Let f be a convex function, defined over the non-empty open convex set X in the topological vector space E. Then f is continuous if, and only if, it is bounded above when restricted to some non-empty open subset U of X.

The condition is obviously necessary, we prove that it is sufficient. Let \(x_{0} \in X\) be a point such that f is bounded above in a neighbourhood V of \(x_{0}\) ; we show firstly that \(f\) is continuous at \(x_0\) . By translation, we can restrict ourselves to the case when \(x_0 = 0\) and \(f(x_0) = 0\) ; moreover we can suppose that the neighbourhood \(V\) is balanced (I, p.~7, prop. 4). Suppose that \(f(x) \leqslant a\) in \(V\) ; for every \(\varepsilon, 0 < \varepsilon < 1\) , we observe that if \(x \in \varepsilon V\) , then \(x / \varepsilon \in V\) and \(-x / \varepsilon \in V\) . Applying inequality (1) of II, p.~16 to the points \(x / \varepsilon\) and 0 and to the number \(\lambda = \varepsilon\) , we see that \(f(x) \leqslant \varepsilon f(x / \varepsilon) \leqslant \varepsilon a\) ; applying it to points \(x\) and \(-x / \varepsilon\) and the number \(\lambda = 1 / (1 + \varepsilon)\) , gives \(f(x) \geqslant -\varepsilon f(-x / \varepsilon) \geqslant -\varepsilon a\) . Thus \(f(x)\) is arbitrarily small in \(\varepsilon V\) , if \(\varepsilon\) is sufficiently small, and \(f\) is continuous at \(x = 0\) .

Now let \(y\) be some point of \(X\) ; since \(X\) is open, there is a number \(\rho > 1\) such that \(z = \rho y\) belongs to \(X\) . Let \(g\) be the homothety \(x \mapsto \lambda x + (1 - \lambda)z\) of centre \(z\) and ratio \(\lambda = 1 - \frac{1}{\rho}\) , which transforms 0 into \(y\) ; for every point \(g(x) \in g(V)\) , we have from (1)

\[
f (g (x)) \leqslant \lambda f (x) + (1 - \lambda) f (z) \leqslant \lambda a + (1 - \lambda) f (z).
\]

Thus \(f\) is bounded above in a neighbourhood of \(y\) and hence, by the first part, is continuous at \(y\) . The proposition is proved.

COROLLARY. --- Every convex function \(f\) defined over an open convex set \(X\) in \(\mathbf{R}^n\) is continuous in \(X\) .

We can suppose that X is not empty. Then there exist, in X, \(n + 1\) affinely independent points \(a_{i}\) ( \(0 \leqslant i \leqslant n\) ) and the convex envelope of these points, S, contains the open non-empty set formed of the points \(\sum_{i=0}^{n} \lambda_{i} a_{i}\) with \(0 < \lambda_{i} < 1\) for all i and \(\sum_{i=0}^{n} \lambda_{i} = 1\) . By II, p.~17, prop. 20, f is bounded above in S and therefore is continuous.

In a topological vector space of infinite dimensions there exist, in general, linear non-continuous forms (II, p.~80, exerc. 25) and thus convex functions that are not continuous at any point.

\subsection{Semi-norms and convex sets}

Let \(\mathbf{E}\) be a vector space over \(\mathbf{R}\) ; a mapping \(p\) of \(\mathbf{E}\) in \(\mathbf{R}\) is positively homogeneous if, for every \(\lambda \geqslant 0\) and all \(x \in \mathbf{E}\) we have

\[
p (\lambda x) = \lambda p (x).\tag{5}
\]

A positively homogeneous function p on E is convex if, and only if, it satisfies axiom \((\mathrm{SN}_{\mathrm{II}})\) of II, p.~1 for all x, y of E;

\[
p (x + y) \leqslant p (x) + p (y).\tag{6}
\]

In fact, if \(p\) is convex, then for \(x, y\) in \(E\) ,

\[
p \left(\frac {1}{2} (x + y)\right) \leqslant \frac {1}{2} p (x) + \frac {1}{2} p (y)
\]

and, by (5), this relation is equivalent to (6). Conversely, if (6) holds, then we also have for all \(\lambda\) such that \(0 < \lambda < 1\) ,

\[
p (\lambda x + (1 - \lambda) y) \leqslant p (\lambda x) + p ((1 - \lambda) y) = \lambda p (x) + (1 - \lambda) p (y)
\]

using (5).

A convex positively homogeneous function on E is called sub-linear.

If \(p\) is a sub-linear function defined on \(\mathbf{E}\) ; then, by II, § 2.8, corollary, for all \(a > 0\) , the set \(\mathrm{V}(p, a)\) (resp. \(\mathrm{W}(p, a)\) ) of points \(x \in \mathrm{E}\) for which \(p(x) \leqslant a\) (resp. \(p(x) < a\) ) is convex; further this set is absorbent, since for all \(x \in \mathrm{E}\) , there exists \(\lambda > 0\) such that \(p(\lambda x) = \lambda p(x) < a\) .

There is a partial converse of this result :

PROPOSITION 22. --- Let A be a convex set, containing 0, in the vector space E. For all \(x \in \mathbf{E}\) , put

\[
p _ {\mathrm{A}} (x) = \inf _ {\rho > 0, x \in \rho \mathrm{A}} \rho\tag{7}
\]

\((0 \leqslant p_{\mathbf{A}}(x) \leqslant \infty)\) . The function \(p_{\mathbf{A}}\) satisfies

\[
p _ {\mathrm{A}} (x + y) \leqslant p _ {\mathrm{A}} (x) + p _ {\mathrm{A}} (y), \quad p _ {\mathrm{A}} (\lambda x) = \lambda p _ {\mathrm{A}} (x)\tag{8}
\]

for all \(x, y\) in \(\mathbf{E}\) and \(\lambda > 0\) . If \(\mathrm{V}(p_{\mathbf{A}}, \alpha)\) (resp. \(\mathrm{W}(p_{\mathbf{A}}, \alpha)\) ) denotes the set of \(x \in \mathbf{E}\) for which \(p_{\mathbf{A}}(x) \leqslant \alpha\) (resp. \(p_{\mathbf{A}}(x) < \alpha\) ), then

\[
\mathrm{W} (p _ {\mathrm{A}}, 1) \subset \mathrm{A} \subset \mathrm{V} (p _ {\mathrm{A}}, 1).\tag{9}
\]

If \(\mathbf{A}\) is absorbent then \(p_{\mathbf{A}}\) is finite (therefore sublinear).

Since the relations \(x \in \rho A\) and \(\lambda x \in \lambda \rho A\) are equivalent when \(\lambda > 0\) , we have \(p_A(\lambda x) = \lambda p_A(x)\) for \(\lambda > 0\) . Let \(x, y\) be two points of \(E\) . If \(x\) (resp. \(y\) ) is not absorbed by \(A\) then \(p_A(x) = +\infty\) (resp. \(p_A(y) = +\infty\) ) and the inequality \(p_A(x + y) \leqslant p_A(x) + p_A(y)\) is obviously true. Suppose there exist \(\alpha > 0\) , \(\beta > 0\) such that \(x \in \alpha A\) , and \(y \in \beta A\) ; then \(x + y \in \alpha A + \beta A = (\alpha + \beta)A\) (II, p.~8, Remark); and thus \(p_A(x + y) \leqslant p_A(x) + p_A(y)\) . The inclusion \(A \subset V(p_A, 1)\) is clearly true. The inclusion \(W(p_A, 1) \subset A\) follows because \(A\) is convex and contains 0. Finally if \(A\) is absorbent then \(p_A\) is obviously finite.

The function \(p_{A}\) defined by (7) is called the gauge of the convex set A. If A is absorbent and symmetric, then \(p_{A}\) is a semi-norm.

PROPOSITION 23. --- Let E be a topological vector space. If A is an open convex set which contains 0, then \(p_{A}\) is finite and continuous, and \(\mathbf{A} = \mathbf{W}(p_{\mathbf{A}}, 1)\) . If A is a closed convex set containing 0, then \(p_{A}\) is lower semi-continuous and \(\mathbf{A} = \mathbf{V}(p_{\mathbf{A}}, 1)\) .

If \(\mathbf{A}\) is open and contains 0, then it is absorbent. For \(x \in \mathbf{A}\) , there exists \(\rho < 1\) such that \(x / \rho \in \mathbf{A}\) , and thus \(p_{\mathbf{A}}(x) < 1\) ; this, combined with (9) gives \(\mathbf{A} = \mathbf{W}(p_{\mathbf{A}}, 1)\) . Since the convex function \(p_{\mathbf{A}}\) is bounded above in the open set \(\mathbf{A}\) , it is continuous in E (II, p.~18, prop. 21).

Suppose A is closed and contains 0. For every \(x \in E\) with \(p_{\mathbf{A}}(x) \leqslant 1\) , we have \(x \in \rho A\) for all \(\rho > 1\) , therefore \(x \in A\) since A is closed; remembering (9), this shows that \(\mathbf{A} = \mathrm{V}(p_{\mathbf{A}}, 1)\) . For all \(\mu > 0\) , \(\mu A\) is therefore the set of \(x\) such that \(p_{\mathbf{A}}(x) \leqslant \mu\) ; as \(p_{\mathbf{A}}(x) \geqslant 0\) in E, this shows that \(p_{\mathbf{A}}\) is lower semi-continuous in E (GT, IV, § 6.2).

A positive sublinear function \(p\) over \(\mathbf{E}\) is the gauge of each convex set \(\mathbf{A}\) where \(\mathrm{W}(p,1)\subset \mathrm{A}\subset \mathrm{V}(p,1)\) .

\section{THE HAHN-BANACH THEOREM (ANALYTIC FORM)}

\subsection{Extension of positive linear forms}

PROPOSITION 1. --- Let E be a preordered vector space and V be a vector subspace of E such that every element of E is bounded above by an element of V. Given a linear form f on V that is positive for the preordered vector space structure of V (induced by that of E) there exists a non-empty set \(S_{f}\) of positive linear forms on E, each being an extension of f.~If \(h \in S_{f}\) then the values \(h(a)\) for \(a \in E\) lie in the interval \([\alpha', \alpha'']\) , where

\[
\alpha^ {\prime} = \sup _ {z \in \mathbf {V}, z \leqslant a} f (z), \quad \alpha^ {\prime \prime} = \inf _ {y \in \mathbf {V}, y \geqslant a} f (y).\tag{1}
\]

I. Special case.

Suppose firstly that \(E = V + R a\) . Since the proposition is trivial if \(a \in V\) , we confine ourselves to the case \(a \notin V\) . The conditions on V imply that the set \(A''\) of \(y \in V\) such that \(a \leqslant y\) is not empty; similarly the set \(A'\) of \(z \in V\) such that \(-z \geqslant -a\) (i.e.~\(z \leqslant a\) ) is not empty. For \(y \in A''\) and \(z \in A'\) , we have \(z \leqslant a \leqslant y\) , and thus by hypothesis \(f(z) \leqslant f(y)\) . Thus \(\alpha', \alpha''\) are finite and \(\alpha' \leqslant \alpha''\) . Any linear form \(f_1\) on E that extends f is completely determined by \(f_1(a)\) and for all \(\lambda \in R\) and all \(x \in V\) , we have

\[
f _ {1} (x + \lambda a) = f (x) + \lambda f _ {1} (a).
\]

Thus \(f_{1}\) is positive if and only if the relations

\[
x \in \mathbf {V}, \quad \lambda \in \mathbf {R}, \quad x + \lambda a \geqslant 0\tag{2}
\]

imply

\[
f (x) + \lambda f _ {1} (a) \geqslant 0.\tag{3}
\]

As \(f(\mu x) = \mu f(x)\) and the relations \(x \geqslant 0\) and \(\mu x \geqslant 0\) are equivalent for \(\mu > 0\) , it is sufficient to show that (2) implies (3) in the particular cases \(\lambda = 0, \lambda = 1\) and \(\lambda = -1\) . For \(\lambda = 0\) , the fact that (2) implies (3) follows from the hypothesis that \(f\) is positive. For \(\lambda = 1\) , to say that (2) implies (3) means that for \(-x \in A'\) , we have \(f_1(a) \geqslant f(-x)\) , i.e.~\(f_1(a) \geqslant \alpha'\) ; for \(\lambda = -1\) , (2) implies (3), means that for \(x \in A''\) , we have \(f(x) \geqslant f_1(a)\) , i.e.~\(f_1(a) \leqslant \alpha''\) . The proposition is therefore proved in this case.

\paragraph{II. General case.}

Let \(\mathfrak{F}\) be the set of pairs \((\mathrm{W}, g)\) where \(\mathrm{W}\) is a vector subspace of \(\mathrm{E}\) containing \(\mathrm{V}\) and \(g\) is a positive linear form on \(\mathrm{W}\) which is an extension of \(f\) . We order \(\mathfrak{F}\) putting \((\mathrm{W}, g) \leqslant (\mathrm{W}', g')\) if \(\mathrm{W} \subset \mathrm{W}'\) and if \(g'\) is an extension of \(g\) . Clearly \(\mathfrak{F}\) is inductive and by th. 2 of S, III, § 2.4, there is a maximal element \((\mathrm{W}_0, g_0)\) . Suppose \(\mathrm{W}_0 \neq \mathrm{E}\) . Then there exists a vector \(b \notin \mathrm{W}_0\) , and, if \(\mathrm{W}_1 = \mathrm{W}_0 + \mathbf{R}b\) , the special case above shows that there exists a positive linear form on \(\mathrm{W}_1\) which is an extension of \(g_0\) ; this contradicts the hypothesis that \((\mathrm{W}_0, g_0)\) is maximal. Thus \(\mathrm{W}_0 = \mathrm{E}\) , and the first part of the proposition is proved. When \(a \in \mathrm{V}\) , the second assertion is obviously true with \(\alpha' = \alpha'' = f(a)\) ; if, on the contrary, \(a \notin \mathrm{V}\) and one puts \(\mathrm{V}_1 = \mathrm{V} + \mathbf{R}a\) , the second assertion follows from the special case I of the proof.

COROLLARY. --- In a topological vector space E with a compatible preorder structure, let P be the set of elements \(\geqslant 0\) in E. Let V be a vector subspace of E containing at least one interior point \(x_{0}\) of P. Then every positive linear form on V can be extended to a positive linear form on E.

By prop. 1 it is sufficient to show that for every \(x \in \mathbf{E}\) , there exists \(x' \in \mathbf{V}\) such that \(x' - x \in \mathbf{P}\) . Now let \(\mathbf{U}\) be a neighbourhood of 0 in \(\mathbf{E}\) such that \(x_0 + \mathbf{U} \subset \mathbf{P}\) . Then \(x + x_0 + \mathbf{U} \subset x + \mathbf{P}\) , and, hence there exists \(\varepsilon\) such that \(0 < \varepsilon < 1\) and the point \(y = x_0 + (1 - \varepsilon)x\) belongs to \(x + \mathbf{P}\) ; then every point of the form \(x + \lambda(y - x)\) belongs to \(x + \mathbf{P}\) for \(\lambda > 0\) . If we take \(\lambda = 1 / \varepsilon\) , then \(x + \lambda(y - x) = \lambda x_0 \in \mathbf{V}\) , from which the conclusion follows.

The conclusion of the corollary is not necessarily valid if one does not assume that V contains an interior point of P, even if E is of finite dimension and if \(P \cap V\) contains points interior in V (II, p.~91, exerc. 25, b)).
